\chapter{The Complete Axiomatization and Exhaustive Search}\label{ch:completeness}

% Local notation of this chapter.
\providecommand{\cmplr}[1]{\texttt{\detokenize{#1}}}
\tikzset{
  cmplbox/.style={rectangle,draw,fill=zxgreen,minimum height=4.2mm,minimum width=4.2mm,
    inner sep=1.5pt,font=\scriptsize},
  cmplpk/.style={circle,draw,fill=zxred!40,minimum size=4.8mm,inner sep=0.3pt,font=\scriptsize},
  cmplpd/.style={circle,draw,fill=zxred!40,minimum size=2.2mm,inner sep=0pt},
  cmplmult/.style={rectangle,draw,fill=white,minimum height=3.8mm,minimum width=4.4mm,
    inner sep=1.5pt,font=\scriptsize},
  cmplmultt/.style={rectangle,draw,fill=gray!20,minimum height=3.8mm,minimum width=4.4mm,
    inner sep=1.5pt,font=\scriptsize},
  cmplwt/.style={isosceles triangle,isosceles triangle apex angle=60,shape border rotate=-90,
    draw,fill=black,minimum size=3mm,inner sep=0pt},
  cmplv/.style={rectangle,draw,fill=white,minimum height=4.2mm,minimum width=5.5mm,
    inner sep=1.5pt,font=\scriptsize},
}
\providecommand{\cmplempty}{\tikz[baseline=-0.5ex]{\draw[dashed,gray] (0,-0.25) rectangle (0.5,0.25);}}
\providecommand{\cmpleq}[3]{\par\smallskip\noindent\hfil#1\quad$=$\quad\ensuremath{#2}\,#3\hfil\par\smallskip}

This chapter specifies the rules that make the rewriting system of ArchytasZX complete for
concrete diagrams, and the exhaustive search that uses them. Two modules implement it.
\module{archytaszx.rewrite.rules_axioms} states every axiom of the qudit ZXW-calculus of Po\'or
et al.~\cite{poor2023zxw} (cited as [P1]) and every axiom of the qufinite ZXW-calculus of Wang,
Po\'or, and Shaikh~\cite{wang2023qufinitezxw} (cited as [P2]) that [P1] lacks, as
\term{template rules}: equations between two small diagrams that the generic template matcher of
\module{archytaszx.rewrite.match_template} finds and rebuilds. It adds the definitional rules that
translate the engine's generators into the papers' generators, templates of the engine's older
hand-written rules so that their converses exist, and the converse of every equation.
\module{archytaszx.rewrite.exhaustive} searches for a derivation between two diagrams by a fair
bidirectional breadth-first enumeration of every application of every registered rule. Together
they realize the transfer of the completeness theorems of [P1] and [P2] to the engine
(\cref{sec:completeness-transfer}): two concrete diagrams with equal denotations are joined by a
finite engine derivation, and exhaustive search finds one.

The template matcher, the \py!Rule! record, the rewrite engine, the comparison view, and the
tactic records \py!ProofPath! and \py!ProofStep! are specified in \cref{ch:rewriting}; the
hand-written rules in \cref{ch:rules}; the generators, phase vectors, and dimensions in
\cref{ch:diagrams,ch:algebra}; the denotation in \cref{ch:semantics}; the command
\azx!prove exhaustive! and the function \py!prove_exhaustive! that wraps the search in
\cref{ch:proofs,ch:shell}. The discard rules of mixed diagrams, which reuse the template builder of
this chapter, are specified in \cref{ch:mixed}.

\section{Sources, Generators, and Conventions}\label{sec:completeness-conventions}

\subsection{The two calculi}

The ZX-calculus~\cite{coecke2011interacting,vandewetering2020working} and the ZW-calculus~\cite{hadzihasanovic2015zw}
combine in the ZXW-calculus. [P1] axiomatizes the ZXW-calculus at one concrete dimension
$d\ge 2$ and proves it complete: two diagrams that denote the same linear map are joined by a
finite sequence of axiom applications in context. [P2] extends the language to wires of
different dimensions, with a dimension splitter and its transpose, building on the qufinite
ZX-calculus~\cite{wang2021qufinite}, and proves this qufinite ZXW-calculus complete. Write
$\mathrm{ZXW}_d$ for the calculus of [P1] at dimension $d$ and $\mathrm{ZXW}_{\mathrm{qf}}$ for the
calculus of [P2]. Neither result covers \bboxes{}~\cite{carette2022bangboxes} or symbolic
dimensions.

The axiom labels of this chapter are those of the papers: S1, S2, D1, Ept, B2, K0, K1, K2, Zer,
P1, H1, Pcy, AD, Sym, Aso, BZW, WW, Bs0, Bsj, TA, HD, ZV, VA, VW, and KZ of [P1], with the
notational equations (Du) and (YT); WF, MW, D1, WW, VA, HD, DD, E, DZX, K0, BZW, TA, and S1 of
[P2], with the equations numbered (2.8), (2.9), and (2.10) and the definition (2.3.3) of the
mixed-dimensional Z box. Rule names follow one scheme:
\begin{itemize}
\item \cmplr{zxw_*} states an axiom of [P1] (several are also axioms of [P2]);
\item \cmplr{qf_*} states an axiom of [P2] that is not one of [P1]'s;
\item \cmplr{*_definition}, \cmplr{phase_to_weight}, \cmplr{weight_to_phase}, and
  \cmplr{scalar_box} are the definitional rules;
\item \cmplr{*_converse} states the same equation read right to left.
\end{itemize}

\subsection{The translation of generators}\label{sec:completeness-translation}

A \term{translation} $\tau$ sends each generator or notation of the papers to an engine diagram
and a scalar factor $s$: the paper generator denotes $s$ times the engine diagram. The module
realizes $\tau$ through the following conventions, fixed in its templates.

\begin{table}[H]
\centering\small
\begin{tabularx}{\linewidth}{@{}>{\raggedright\arraybackslash}p{0.33\linewidth}>{\raggedright\arraybackslash}X>{\raggedright\arraybackslash}p{0.17\linewidth}@{}}
\toprule
Paper generator or notation & Engine diagram & Factor $s$ \\
\midrule
Z box with parameter $a$ & Z spider whose phase vector is the weight vector with coefficients $a$ & $1$ \\
green circle spider with phases $\alpha$ & Z spider with phase vector $\alpha$ & $1$ \\
Hadamard box $H$; $H^\dagger$ & $F$; $F^3$ (three $F$ in series) & $1$ \\
$D$ & $F^2$ & $1$ \\
pink spider $K_j$, $m$ inputs, $n$ outputs & X spider with phase vector $K_j$ & $d^{(m+n)/2-1}$ \\
pink dot (phaseless pink spider) & phaseless X spider & $d^{(m+n)/2-1}$ \\
plain green dot & phaseless Z spider & $1$ \\
W node; W transpose & W node; W node with its ports bent & $1$ \\
cap, cup, swap, identity & wires & $1$ \\
triangle; inverse triangle & $T$; $T_i$ & $1$ \\
dimension splitter; merger & $S$; $B$ & $1$ \\
closed scalar diagram & its value, in the diagram scalar & $1$ \\
multiplier, V box, M box & their definitions (\cref{sec:completeness-notations}) & as built \\
mixed-dimensional Z box & its definition (2.3.3) (\cref{sec:completeness-mixed}) & $1$ \\
\bottomrule
\end{tabularx}
\caption{The translation $\tau$ that the templates of \module{archytaszx.rewrite.rules_axioms} encode.}
\label{tab:completeness-translation}
\end{table}

A W transpose is a W node whose input and outputs are bent: its outputs are wired to output
ports elsewhere or its input to an input port elsewhere. A cap or cup is a wire between two ports
of the same direction. Neither needs a generator of its own.

\begin{definition}[Pink normalization]\label{def:completeness-pink}
A \term{pink node} is a template node flagged \py!pink=True!: an X spider in [P1]'s
normalization. A pink node with $k$ ports at dimension $D$ denotes
\[
  \nu_D(k) \;=\; D^{k/2-1}
\]
times the engine's X spider with the same phase vector. $D$ is the node's own dimension when the
template gives one (\py!TNode.dim!), else the dimension of the host match. The number $k$ counts
the template ports plus, for an open node, the host ports it carries beyond them.
\end{definition}

In the pictures of this chapter a light-red node is a pink node and a full-red node an engine X
spider; a green circle is a Z spider with a phase vector, a green square a Z box (a Z spider
read through its coefficients), a small green dot a plain green dot, a small light-red dot a pink
dot, a black upward triangle a W node, a black downward triangle a W transpose, a yellow square
the Fourier box $F$, and a white triangle the triangle $T$. Diagrams flow from bottom to top.

\subsection{Coefficient vectors}\label{sec:completeness-vectors}

A Z spider at concrete dimension $d$ carries a phase vector $v$ (\cref{ch:algebra}); its
\term{coefficients} are $c_0=1$ and $c_k=w_k\,e^{2\pi i t_k}$ for $1\le k\le d-1$, where $t_k$ is the
phase in turns at index $k$ and $w_k$ the weight, both read from the stored entries or rules of
$v$ (\py!PhaseVector.coefficient(k)!). A phaseless spider (\py!None!) has all coefficients $1$.
The Z box with parameter $a=(a_1,\dots,a_{d-1})$ denotes
$\sum_{k=0}^{d-1} a_k\,\ket{k\cdots k}\bra{k\cdots k}$ with $a_0=1$. Indices of $a$ are read
modulo $d$, so $a_{-j}=a_{d-j}$. Addition of phase vectors (\py!PhaseVector.__add__!) adds
phases and multiplies weights, hence multiplies coefficients; it is the product $ab$ of Z boxes
with $(ab)_k=a_kb_k$.

The module defines these vector operations. Each takes a \py!Dim! and returns a
\py!PhaseVector! or \py!None! (phaseless).

\begin{signature}
\py!coefficients(pv: PhaseVector | None, dim: Dim) -> list[Scalar]!\\
$[c_0,\dots,c_{d-1}]$ at the concrete dimension \py!dim! (\py!dim.to_int()! raises at a symbolic
one); \py!None! gives $d$ copies of $1$.

\py!from_coefficients(values: Sequence[Scalar], dim: Dim) -> PhaseVector | None!\\
The weight vector with weights \py!values[1:]! at indices $1,\dots,\text{len}-1$
(\py!PhaseVector.from_weights!); \py!values[0]! is ignored and taken as $1$. Returns \py!None!
when the result \py!is_zero! (the plain vector).

\py!vector_of(values, dim)! $=$ \py!from_coefficients([1, *values], dim)!: the Z box
$(\text{values}[0],\dots,\text{values}[d-2])$.
\end{signature}

\begin{signature}
\py!zero_vector(dim)!: weights all $0$, the Z box $(0,\dots,0)=\ket{0}\bra{0}$ (a weight rule $0$).\\
\py!minus_ones(dim)!: weights all $-1$ (a weight rule $-1$).\\
\py!constant_vector(x: Scalar, dim)!: weights all $x$ (a weight rule $x$).\\
\py!last_vector(x: Scalar, dim)!: the \term{box labelled $x$}, $(0,\dots,0,x)$: weight $0$ at
$1,\dots,d-2$ and $x$ at $d-1$. At a symbolic dimension it raises \py!RewriteDomainError!
\code!"the box (0, ..., 0, x) needs a concrete dimension, got {dim}"!.\\
\py!unit_vector(index, dim)!: at concrete $d$, with $i=\text{index}\bmod d$, weight $1$ at $i$ when
$i\ne 0$ and $0$ at every other index of $1,\dots,d-1$; write $u_i$ for it ($u_0=(0,\dots,0)$). At a symbolic
dimension, the weight rule \py!KroneckerDelta(INDEX, 0 + index)!.
\end{signature}

\begin{signature}
\py!reversed_vector(pv, dim)!: $\overleftarrow{a}=(a_{d-1},\dots,a_1)$, the coefficient at $k$ being
$a_{d-k}$; \py!None! and a zero vector pass through unchanged. An unweighted vector keeps its phase
form (entry at $k$ is the entry of \py!pv! at $d-k$); a weighted one is rebuilt with
\py!from_coefficients!. Concrete dimension only.

\py!sum_vector(a, b, dim) -> PhaseVector!: $a+b$ entrywise on indices $1,\dots,d-1$, as weights.
At a symbolic dimension both vectors must carry neither explicit entries nor explicit weights
(else \py!RewriteDomainError! \code!"a + b at a symbolic dimension needs rule-only vectors"!); the
result is the weight rule $r_a+r_b$ with $r_v=w\,e^{2\pi i t}$, $w$ the weight rule of $v$ (or $1$)
and $t$ its phase rule in turns (or absent).

\py!shifted_vector(pv, j: int, dim)!: $k_j(a)=(a_{1-j}/a_{-j},\dots,a_{d-1-j}/a_{-j})$, the
coefficient at $k$ being $a_{k-j}/a_{-j}$. It calls the memoized
\py!_shifted_vector(pv, j mod d, dim)! (an \py!lru_cache! of 256 entries), which computes
$1/a_{-j}$ once and raises \py!RewriteDomainError! \code!"k_j(a) needs a_{-j} != 0"! when the
pivot \py!is_zero! or simplifies to $0$.
\end{signature}

Two predicates decide zero-ness. \py!_nonzero(x)! holds when \py!x.is_zero! is false and
\py!sympy.simplify! of $x$ is not $0$. \py!_zero(expr)! first evaluates a symbol-free
expression to 30 digits and returns false when the modulus exceeds $10^{-8}$; otherwise (or when
evaluation raises \py!TypeError!, \py!ValueError!, or \py!ArithmeticError!) it returns whether
\py!simplify(expr) == 0!. \py!_same_coefficients(x, y, dim)! holds when the vectors are equal as
objects (\py!None! read as the plain vector), or, at a concrete dimension, when $\_zero(x_k-y_k)$
holds for every $k=1,\dots,d-1$.

\subsection{The phase vector \texorpdfstring{$K_j$}{K\_j}}

The phase vector $K_j$ has phase $jk/d$ turns at index $k$ (function \py!k_vector! of
\module{archytaszx.rewrite.match_template}). At a concrete dimension $j$ must be an integer;
$j\bmod d=0$ gives \py!None! (phaseless), otherwise the memoized vector with entries
$(jk\bmod d)/d$. At a symbolic dimension $K_j$ is the phase rule $j\cdot\mathtt{INDEX}/d$, and
$j=0$ gives \py!None!. The pink spider $K_j$ of [P1] is the pink node with phase vector $K_j$;
the pink dot is $K_0$.

\subsection{Composite notations}\label{sec:completeness-notations}

\begin{definition}[Multiplier]\label{def:completeness-multiplier}
The \term{multiplier} labeled $m\ge 0$ is a plain green dot with one input and $m$ outputs whose
outputs are the $m$ inputs of a pink dot with one output (\py!TB.multiplier(m)!). Its
\term{transpose} (\py!up=True!) is a pink dot with one input and $m$ outputs into a plain green
dot with $m$ inputs and one output. Both have one input and one output and carry the pink
normalization $\nu_d(m+1)=d^{(m-1)/2}$. \py!_mult_at(t, m, dim, up=...)! builds the same shape
with every node of the explicit dimension \py!dim!.
\end{definition}

\begin{center}
\begin{tikzpicture}[zxd]
\node[cmplmult] (m) at (0,0) {$m$};
\draw[wire] (0,-0.7) -- (m) -- (0,0.7);
\node at (0.9,0) {$:=$};
\begin{scope}[xshift=2cm]
\node[zdot] (g) at (0,-0.4) {};
\node[cmplpd] (p) at (0,0.4) {};
\draw[wire] (0,-0.9) -- (g); \draw[wire] (p) -- (0,0.9);
\draw[wire] (g) to[bend left=50] (p); \draw[wire] (g) to[bend right=50] (p);
\node[dimlabel] at (0,0) {$\cdots$}; \node[dimlabel] at (0.55,0) {$m$};
\end{scope}
\begin{scope}[xshift=4.6cm]
\node[cmplmultt] (m) at (0,0) {$m$};
\draw[wire] (0,-0.7) -- (m) -- (0,0.7);
\node at (0.9,0) {$:=$};
\end{scope}
\begin{scope}[xshift=6.6cm]
\node[cmplpd] (g) at (0,-0.4) {};
\node[zdot] (p) at (0,0.4) {};
\draw[wire] (0,-0.9) -- (g); \draw[wire] (p) -- (0,0.9);
\draw[wire] (g) to[bend left=50] (p); \draw[wire] (g) to[bend right=50] (p);
\node[dimlabel] at (0,0) {$\cdots$}; \node[dimlabel] at (0.55,0) {$m$};
\end{scope}
\end{tikzpicture}
\end{center}

\noindent A white box labeled $m$ is a multiplier; a gray box is a transposed multiplier. A
multiplier labeled $0$ is a plain green effect followed by a pink dot state.

\begin{definition}[V box and M box]\label{def:completeness-vm}
At concrete dimension $d$, the \term{V box} (\py!_v_box(t, d)!) is the box labeled $1$,
$(0,\dots,0,1)$, followed by a W node with $d-1$ outputs; output $k-1$ ($1\le k\le d-1$) passes
through the multiplier labeled $(-k)\bmod d$ into input $k-1$ of a pink dot with $d-1$ inputs
and one output. The \term{M box} (\py!_m_box(t, d)!) is the box labeled $1$ followed by a W
node with $d$ outputs; output $p$ passes through the multiplier labeled $L_p$ into input $p$ of
a pink dot with $d$ inputs, where $(L_0,\dots,L_{d-1})=((-1)\bmod d,\dots,(-(d-1))\bmod d,0)$.
Both have one input and one output.
\end{definition}

\begin{center}
\begin{tikzpicture}[zxd]
\node[cmplv] (v) at (0,0) {$V$}; \draw[wire] (0,-0.8) -- (v) -- (0,0.8);
\node at (0.9,0) {$:=$};
\begin{scope}[xshift=2.6cm]
\node[cmplbox] (b) at (0,-1.3) {$1$};
\node[wnode] (w) at (0,-0.55) {};
\node[cmplmult] (m1) at (-0.8,0.35) {$-1$};
\node[cmplmult] (m2) at (0.8,0.35) {$1$};
\node[dimlabel] at (0,0.35) {$\cdots$};
\node[cmplpd] (p) at (0,1.2) {};
\draw[wire] (0,-1.9) -- (b) -- (w);
\draw[wire] (w) -- (m1) -- (p); \draw[wire] (w) -- (m2) -- (p); \draw[wire] (p) -- (0,1.8);
\end{scope}
\begin{scope}[xshift=5.6cm]
\node[cmplv] (v) at (0,0) {$M$}; \draw[wire] (0,-0.8) -- (v) -- (0,0.8);
\node at (0.9,0) {$:=$};
\end{scope}
\begin{scope}[xshift=8.4cm]
\node[cmplbox] (b) at (0,-1.3) {$1$};
\node[wnode] (w) at (0,-0.55) {};
\node[cmplmult] (m1) at (-1.1,0.35) {$-1$};
\node[cmplmult] (m2) at (0.2,0.35) {$1$};
\node[cmplmult] (m3) at (1.1,0.35) {$0$};
\node[dimlabel] at (-0.45,0.35) {$\cdots$};
\node[cmplpd] (p) at (0,1.2) {};
\draw[wire] (0,-1.9) -- (b) -- (w);
\draw[wire] (w) -- (m1) -- (p); \draw[wire] (w) -- (m2) -- (p); \draw[wire] (w) -- (m3) -- (p);
\draw[wire] (p) -- (0,1.8);
\end{scope}
\end{tikzpicture}
\end{center}

\noindent The labels $-1,\dots,1$ stand for $(-1)\bmod d=d-1$ down to $(-(d-1))\bmod d=1$. A V box carries the pink
normalization $\nu_d(d)\prod_{L=1}^{d-1}\nu_d(L+1)=d^{\,d/2-1+(d-1)(d-2)/4}$; an M box carries
$d^{(d-2)(d+1)/4}$.

\begin{definition}[Mixed-dimensional Z box]\label{def:completeness-mixedz}
For ports $(\delta_1,d_1),\dots,(\delta_r,d_r)$ (direction $\delta_i$, concrete dimension
$d_i$), let $D=\prod_i d_i$ and $\ell=\min_i d_i$. The \term{mixed-dimensional Z box}
(\py!mixed_z(t, legs, param)!, definition (2.3.3) of [P2]) is a Z spider of dimension $D$ with one
input per input port and one output per output port. An input port $i$ is port $1$ of a merger
$B$ with input dimensions $(D/d_i,d_i)$ and output dimension $D$; port $0$ of that $B$ carries a
pink dot state of dimension $D/d_i$, and its output feeds the Z spider. An output port $i$ is
port $1$ of a splitter $S$ from $D$ to $(D/d_i,d_i)$ fed by the Z spider; port $0$ of that $S$
ends in a pink dot effect of dimension $D/d_i$. The Z spider's parameter is \py!param!, by default
the fixed vector with coefficients $1$ below $\ell$ and $0$ from $\ell$ on (\py!_ones_below!); this
default is the \term{plain mixed green dot}. The function returns the external ports in the order of
\py!legs!. Its pink normalization is $\prod_i (D/d_i)^{-1/2}$.
\end{definition}

\begin{center}
\begin{tikzpicture}[zxd]
\node[cmplbox] (z) at (0,0) {$a$};
\draw[wire] (z) -- (-0.4,-0.8) node[below,dimlabel] {$d_1$};
\draw[wire] (z) -- (0.4,-0.8) node[below,dimlabel] {$d_2$};
\draw[wire] (z) -- (0,0.8) node[above,dimlabel] {$d_3$};
\node at (1.1,0) {$:=$};
\begin{scope}[xshift=3.4cm]
\node[zs,minimum size=5mm] (z) at (0,0) {$a$};
\node[bnode] (b1) at (-0.9,-1) {$B$};
\node[bnode] (b2) at (0.9,-1) {$B$};
\node[snode] (s) at (0,1) {$S$};
\node[cmplpd] (p1) at (-1.25,-1.75) {};
\node[cmplpd] (p2) at (0.55,-1.75) {};
\node[cmplpd] (p3) at (-0.35,1.75) {};
\draw[wire] (p1) -- (b1.south west); \draw[wire] (b1.south east) -- (-0.55,-2.1) node[below,dimlabel] {$d_1$};
\draw[wire] (p2) -- (b2.south west); \draw[wire] (b2.south east) -- (1.25,-2.1) node[below,dimlabel] {$d_2$};
\draw[wire] (b1) -- (z); \draw[wire] (b2) -- (z); \draw[wire] (z) -- node[right,dimlabel] {$D$} (s);
\draw[wire] (s.north west) -- (p3); \draw[wire] (s.north east) -- (0.35,2.1) node[above,dimlabel] {$d_3$};
\end{scope}
\end{tikzpicture}
\end{center}

\subsection{The scalar convention}\label{sec:completeness-scalar}

A template rule states $L = s\cdot R$, where $L$ and $R$ are the two templates in paper
normalization. The stated factor is assembled from three parts of the rule: the axiom's own
\py!scalar(dim, bindings)!, and the \py!scalar! callables the two templates carry (a scalar
diagram of the paper on that side, already evaluated). In the pictures below, a template scalar is
written as a factor in front of its side and the axiom scalar in front of the right side. Applying the
rule left to right replaces an occurrence of the engine image of $L$ by the engine image of $R$
and multiplies the diagram scalar by the \term{engine scalar} (function \py!template_scalar!)
\begin{equation}\label{eq:completeness-engine-scalar}
  \sigma \;=\; s_{\mathrm{ax}}\cdot\frac{N_R}{N_L}\cdot\frac{s_R}{s_L},
  \qquad N_X=\prod_{\text{pink }v\in X}\nu_{D_v}(k_v),
\end{equation}
where $s_{\mathrm{ax}}$ is the axiom scalar, $s_L$ and $s_R$ are the template scalars (absent means
$1$), and $k_v$ is the port count of the pink node $v$ in the match, including the host ports an
open node carries. \py!_inverse! raises \py!RewriteDomainError!
\code!"a template rule scalar of 0 has no converse"! when it must invert $0$. When both sides
have one shape at every dimension, carry no template scalar, and have no open node,
\py!fixed_template_scalar! of \module{archytaszx.rewrite.rules_library} evaluates $\sigma$ once at
the symbol $d$ and records it as the rule's \py!scalar_introduced!; otherwise
\py!scalar_introduced! is $1$ and \py!scalar_in_match! computes $\sigma$ per match. The match is
rejected when $s_L$ evaluates to $0$.

\section{Template Machinery of the Module}\label{sec:completeness-machinery}

\subsection{Parameters}

A template node carries a \py!Param! (\cref{ch:rewriting}): \py!match(phase, dim, bindings)!
extends the bindings by the match of a host phase vector or returns \py!None!;
\py!build(dim, bindings)! returns the vector the parameter denotes; \py!variables()! names the
variables \py!build! reads; \py!binds()! names those \py!match! binds. The matcher module
supplies \py!Fixed! (a function of the dimension; \py!PHASELESS! is \py!Fixed(lambda d: None)!),
\py!Var(name)! (binds the host vector, the plain vector for a phaseless host, or compares it with
the bound value), \py!Derived(reads, compute)! (recomputed from the bindings and compared;
binds nothing), and \py!Kj(name, sign=1, offset=0)! (recognizes $K_{\mathit{sign}\cdot j+\mathit{offset}}$ and binds
$j$, reduced modulo $d$ at a concrete dimension). This module adds three parameters and four
fixed vectors.

\begin{signature}
\py!class LastEntry(Param)!, \py!LastEntry(name: str)!\\
Matches the box labeled $x$. At a symbolic dimension it fails. At concrete $d$ every
coefficient at $1,\dots,d-2$ must be $0$ (\py!is_zero! or simplifying to $0$); $x=c_{d-1}$ is bound
to \py!name!, or, when \py!name! is bound, compared by simplifying the difference to $0$. Builds
\py!last_vector(x, dim)!. \py!variables()! and \py!binds()! are \py!{name}!. Equality and
hash are by \py!name!.
\end{signature}

\begin{signature}
\py!class MixedVar(Param)!, \py!MixedVar(name: str, low: int)!\\
The parameter of a mixed-dimensional Z box whose least port dimension is \py!low!. Matching
requires every coefficient at an index $\ge$ \py!low! to be exactly zero (\py!is_zero!) and binds
the vector of dimension \py!low! with coefficients $c_1,\dots,c_{\mathit{low}-1}$ (the plain vector
when that is zero); a bound value is compared by \py!_same_coefficients! at dimension \py!low!.
Builds \py!_spread(a, low, dim)!: weights $a_k$ for $k<\mathit{low}$ and $0$ for $k\ge\mathit{low}$
over \py!dim!.
\end{signature}

\begin{signature}
\py!class UnitIndex(Param)!, \py!UnitIndex(name: str, sign: int = 1)!\\
Matches $u_{\mathit{sign}\cdot j}$ and binds $j$. At a symbolic dimension it fails. At concrete $d$
every coefficient at $1,\dots,d-1$ must be $0$ (\py!is_zero!) or simplify to $1$, with at most one
$1$; the index $i$ of that $1$ (or $0$) gives the binding $(\mathit{sign}\cdot i)\bmod d$, which
replaces any previous binding of \py!name!. Builds \py!unit_vector(sign * j, dim)!.
\end{signature}

\begin{signature}
\py!ZERO = Fixed(zero_vector)!, \py!MINUS_ONES = Fixed(minus_ones)!,
\py!ONE_LAST = Fixed(lambda dim: last_vector(1, dim))! (the box labeled $1$),
\py!k_fixed(j) = Fixed(lambda dim: k_vector(j, dim))!, \py!fixed(build) = Fixed(build)!.
\end{signature}

\subsection{The template builder}

\begin{signature}
\py!class TB! --- builds a \py!Template! node by node. Fields: \py!nodes: list[TNode]!,
\py!wires: list[tuple[TPort, TPort]]!, \py!terminals: dict[str, TPort]!,
\py!bare: list[tuple[str, str]]!, \py!rest: dict[str, str]!, and a counter for generated keys.
\begin{itemize}[leftmargin=1.2em]
\item \py!node(gen, ins, outs, param=PHASELESS, key=None, *, open=False, pink=False, dim=None) -> str!
  appends a \py!TNode! and returns its key; a missing key becomes \code!_1!, \code!_2!, \ldots{}
  in creation order. After each node and each wire it calls \py!check_template_budget! with the
  node count plus the wire count, which raises \py!TemplateTooLarge! when the template outgrows the host
  being matched.
\item \py!z(ins, outs, param, key, *, open)!, \py!pink(ins, outs, param, key)! (an X spider with
  \py!pink=True!), \py!w(outs, key)! (a W node with one input), \py!f(key)! (a Fourier box).
\item \py!conn(gen, dims, key)! a merger or splitter with explicit input and output dimensions
  \py!dims = ((d_in, ...), (d_out, ...))!, stored as \py!port_dims!.
\item \py!wire(a, b)!, \py!term(name, port)!, \py!chain(keys)! (output $0$ of each node into input
  $0$ of the next), \py!fs(count)! (\py!count! Fourier boxes in series),
  \py!multiplier(m, *, up=False) -> (input, output)!.
\item \py!done(scalar=None) -> Template! freezes the template with the optional template scalar.
\end{itemize}
\py!i(key, index=0)! and \py!o(key, index=0)! name the input and the output port \py!index! of node
\py!key!; \py!_cdim(v)! is \py!Dim.concrete(v)!.
\end{signature}

A template's first node is its \term{anchor}: the matcher tries every host node of the anchor's
generator. An \term{open} node may carry host ports beyond its template ports in its symmetric
groups (\cref{ch:rewriting}); its template \py!rest! map names, per open node of the matched side,
the node of the built side that receives those ports. A \py!bare! pair joins two terminals by a
bare wire on the built side.

\subsection{Sides and arity enumerators}

A \py!Side! couples a template function \py!template(dim, arities)! with an arity enumerator
\py!arities(diagram, anchor, dim)! and a lower bound \py!min_size(dim)! on the template's node plus
wire count; the matcher skips a side whose bound exceeds the host's node plus wire count.
\py!fixed_side(build)! makes a side whose template ignores the dimension and the arities (one
empty arity assignment). \py!concrete_side(build, grows=False)! makes a side whose template
\py!build(d)! exists only at a concrete dimension (\py!None! otherwise); with \py!grows=True! its
\py!min_size! is $d$. \py!arity_from_anchor(build)! lifts an enumerator that ignores the dimension.
\py!_first_kind(gen)! yields one empty assignment when the anchor has generator \py!gen!.

The enumerators read the host through \py!_partners(diagram)! (each wire end to the other),
\py!_legs(diagram, n)! (the input port references, then the output port references, of node
$n$), \py!_neighbour_ids(diagram, n)! (the node at the far end of each connected port of $n$,
in port order, with repetition), and \py!_count(diagram, n)! (the port count of $n$). For the
mixed rules, \py!_host_values(diagram)! is the sorted set of every concrete port dimension of the
diagram, read through its parameter environment (\py!Diagram.resolve_dim!), and
\py!_divisor_pairs(values)! yields every $(m,n)$ with $m\ge 2$, $m\in\text{values}$, and $mn\in\text{values}$
for some $n\ge2$, in the order of \py!values! for $m$ and then for $mn$.

An arity enumerator bounds the instances of a family that a host can match: an instance whose
arity assignment the enumerator does not yield is never tried. The enumerator of each rule is
given with the rule.

\subsection{Template axioms}

A \py!TemplateAxiom! (\cref{ch:rewriting}) has the fields \py!name!, \py!left!, \py!right!,
\py!scalar!, \py!conditions! (named predicates on the bindings, checked on the matched side),
\py!concrete_only! (match only where the anchor's dimension is concrete), \py!description!,
\py!quantifiers!, \py!free!, \py!converse_free!, \py!solve!, and \py!candidates!. A \term{free
variable} is one the built side reads and the matched side does not bind; \py!free! maps it to its
default value, which every match carries, and \azx!apply ... with NAME=VALUE! supplies another,
which the conditions then check. \py!converse_free! holds the free variables of the converse.
\py!solve(dim, bindings)! returns the free variables that the conditions determine from the
others, or \py!None!. \py!candidates(dim, bindings)! lists further values of the free variables,
computed from the matched side; exhaustive search enumerates a match with each
(\py!TemplatePattern.candidate_matches!).

\begin{definition}[Converse]\label{def:completeness-converse}
\py!axiom.converse(name)! is the template axiom with the sides swapped, the scalar
$1/s_{\mathrm{ax}}$ (through \py!_inverse!), the same conditions, \py!concrete_only!, and
quantifiers, the description followed by \code!" Read right to left."!, and \py!free! set to the
original's \py!converse_free!. It carries no \py!converse_free!, \py!solve!, or \py!candidates!.
The template scalars swap sides with the templates, so \cref{eq:completeness-engine-scalar}
gives the reciprocal engine scalar.
\end{definition}

\py!template_rule(axiom)! of \module{archytaszx.rewrite.rules_library} turns a template axiom into
a \py!Rule! with the axiom's name, a \py!TemplatePattern!, the template builder, the axiom's side
conditions, the quantifiers (leg counts: every quantifier other than \code!d!; dimensions:
\code!("d",)!), and \py!automatic! true only for the names in \py!PROMOTED_RULE_NAMES!:
\cmplr{zxw_ept}, \cmplr{scalar_box}, \cmplr{zxw_k0}, \cmplr{zxw_k2}, \cmplr{zxw_zer},
\cmplr{zxw_ad}, \cmplr{zxw_bs0}, \cmplr{zxw_bsj}, \cmplr{zxw_ta}, \cmplr{zxw_bzw},
\cmplr{zxw_aso}, \cmplr{zxw_hd}, \cmplr{zxw_va}, \cmplr{zxw_zv}, \cmplr{zxw_kz}, \cmplr{qf_dzx},
\cmplr{qf_e}, \cmplr{qf_k0}, \cmplr{qf_bzw}, and \cmplr{qf_s1}. Every other template rule,
converses included, is manual: only an explicit \azx!apply!, the tactics, and exhaustive search
use it.

\providecommand{\cmplhead}[2]{\par\medskip\noindent\textbf{\cmplr{#1}}\quad #2\par\nopagebreak\smallskip}

\section{The Axioms of [P1]}\label{sec:completeness-p1}

This section states the 27 rules that encode the axioms of [P1], in the order of
\py!AXIOMS!. Each entry gives the equation as a picture with the stated factors, the engine scalar
$\sigma$ of \cref{eq:completeness-engine-scalar}, the side conditions beyond the template match,
whether the rule is \py!concrete_only!, the arity enumerators of the two sides (the left side's
enumerator serves the forward rule, the right side's the converse), and the converse. A
\term{fixed side} has the single empty arity assignment; its anchor is still required to have the
generator of the template's first node.

\cmplhead{zxw_s1}{[P1] S1: spider fusion of Z boxes.}
\cmpleq{\begin{tikzpicture}[zxd]
\node[cmplbox] (a) at (0,0) {$a$};
\node[cmplbox] (b) at (0,1.4) {$b$};
\draw[wire] (a) to[bend left=45] (b); \draw[wire] (a) to[bend right=45] (b);
\node[dimlabel] at (0,0.7) {$\cdots$}; \node[dimlabel] at (0.62,0.7) {$k$};
\draw[wire] (a) -- (-0.45,-0.7); \draw[wire] (a) -- (0.45,-0.7); \node[dimlabel] at (0,-0.6) {$\cdots$};
\draw[wire] (b) -- (-0.45,2.1); \draw[wire] (b) -- (0.45,2.1); \node[dimlabel] at (0,2.0) {$\cdots$};
\end{tikzpicture}}{}{\begin{tikzpicture}[zxd]
\node[cmplbox] (c) at (0,0.7) {$ab$};
\draw[wire] (c) -- (-0.45,-0.7); \draw[wire] (c) -- (0.45,-0.7); \node[dimlabel] at (0,-0.6) {$\cdots$};
\draw[wire] (c) -- (-0.45,2.1); \draw[wire] (c) -- (0.45,2.1); \node[dimlabel] at (0,2.0) {$\cdots$};
\end{tikzpicture}}
The left template is a Z spider \code!a! with no inputs and $k$ outputs, all wired to the $k$
inputs of a Z spider \code!b! with no outputs; both are open, with \py!Var("a")! and
\py!Var("b")!. The right template is one Z spider \code!c! with no template ports whose parameter
is the derived product $ab$ (\py!PhaseVector! addition), and \code!c! receives the remaining
ports of both (\py!rest = {"a": "c", "b": "c"}!). $\sigma=1$. Left arities: when the anchor is a
Z spider, count, for each other Z spider adjacent to it, the wires joining the two; yield
$\{k\}$ for each distinct count in increasing order. The right side is matched only through the
explicit converse \cmplr{zxw_s1_converse} (\cref{sec:completeness-converses}). Quantifiers
\code!k!, \code!legs!, \code!d!.

\cmplhead{zxw_s2}{[P1] S2: the plain green dot with two ports is the identity.}
\cmpleq{\begin{tikzpicture}[zxd]
\node[zdot] (z) at (0,0) {}; \draw[wire] (0,-0.7) -- (z) -- (0,0.7);
\end{tikzpicture}}{}{\begin{tikzpicture}[zxd]\draw[wire] (0,-0.7) -- (0,0.7);\end{tikzpicture}}
Left: a phaseless Z spider with one input (terminal \code!in!) and one output (terminal
\code!out!). Right: the bare wire \code!(in, out)!. $\sigma=1$; fixed sides. The generated
converse \cmplr{zxw_s2_converse} inserts a plain green dot on a wire: one match per host wire,
and one per boundary port whose node is not a merger or splitter (\cref{ch:rewriting}).

\cmplhead{zxw_s2_pink}{[P1] S2 (second part), [P2] S2: the pink dot with two ports is the identity.}
\cmpleq{\begin{tikzpicture}[zxd]
\node[cmplpd] (z) at (0,0) {}; \draw[wire] (0,-0.7) -- (z) -- (0,0.7);
\end{tikzpicture}}{}{\begin{tikzpicture}[zxd]\draw[wire] (0,-0.7) -- (0,0.7);\end{tikzpicture}}
Left: a pink dot with one input and one output; right: the bare wire. $\nu_d(2)=1$, so
$\sigma=1$. Fixed sides; generated converse \cmplr{zxw_s2_pink_converse}.

\cmplhead{zxw_d1}{[P1] D1: a Z box passes through $D$ reversed.}
\cmpleq{\begin{tikzpicture}[zxd]
\node[cmplbox] (a) at (0,0) {$a$}; \node[fbox] (f1) at (0,0.75) {$F$}; \node[fbox] (f2) at (0,1.45) {$F$};
\draw[wire] (0,-0.6) -- (a) -- (f1) -- (f2) -- (0,2.05);
\end{tikzpicture}}{}{\begin{tikzpicture}[zxd]
\node[fbox] (f1) at (0,0) {$F$}; \node[fbox] (f2) at (0,0.7) {$F$}; \node[cmplbox] (a) at (0,1.45) {$\overleftarrow{a}$};
\draw[wire] (0,-0.6) -- (f1) -- (f2) -- (a) -- (0,2.05);
\end{tikzpicture}}
Left: a Z spider (one input, one output) with \py!Var("a")! followed by two $F$ in series. Right:
two $F$ followed by a Z spider with the derived parameter \py!reversed_vector(a)!. $\sigma=1$.
\py!concrete_only!. Converse: \cmplr{zxw_d1_converse}.

\cmplhead{zxw_ept}{[P1] Ept: the pink dot state into a Z box effect is the empty diagram.}
\cmpleq{\begin{tikzpicture}[zxd]
\node[cmplpd] (p) at (0,0) {}; \node[cmplbox] (a) at (0,0.8) {$a$}; \draw[wire] (p) -- (a);
\end{tikzpicture}}{}{\cmplempty}
Left: a pink dot with one output into a Z spider with one input and \py!Var("a")!. Right: the
empty template. $\sigma=\nu_d(1)^{-1}=\sqrt d$; fixed sides; promoted to automatic.
\py!converse_free = {"a": plain vector}!, so the generated converse \cmplr{zxw_ept_converse}
inserts a pink dot state into a plain green effect by default and accepts any $a$ through
\azx!with a=...!. A side with no nodes and no bare wire matches once per distinct dimension of
the host's nodes (\cref{ch:rewriting}).

\cmplhead{zxw_b2}{[P1] B2, [P2] B2: the bialgebra of the plain green and pink dots.}
\cmpleq{\begin{tikzpicture}[zxd]
\node[zdot] (g1) at (-0.5,0) {}; \node[zdot] (g2) at (0.5,0) {};
\node[cmplpd] (p1) at (-0.5,1.1) {}; \node[cmplpd] (p2) at (0.5,1.1) {};
\draw[wire] (g1) -- (p1); \draw[wire] (g1) -- (p2); \draw[wire] (g2) -- (p1); \draw[wire] (g2) -- (p2);
\draw[wire] (g1) -- (-0.5,-0.6) node[below,dimlabel] {$a$}; \draw[wire] (g2) -- (0.5,-0.6) node[below,dimlabel] {$b$};
\draw[wire] (p1) -- (-0.5,1.7) node[above,dimlabel] {$c$}; \draw[wire] (p2) -- (0.5,1.7) node[above,dimlabel] {$e$};
\end{tikzpicture}}{}{\begin{tikzpicture}[zxd]
\node[cmplpd] (p) at (0,0.2) {}; \node[zdot] (g) at (0,0.9) {}; \draw[wire] (p) -- (g);
\draw[wire] (p) -- (-0.4,-0.6) node[below,dimlabel] {$a$}; \draw[wire] (p) -- (0.4,-0.6) node[below,dimlabel] {$b$};
\draw[wire] (g) -- (-0.4,1.7) node[above,dimlabel] {$c$}; \draw[wire] (g) -- (0.4,1.7) node[above,dimlabel] {$e$};
\end{tikzpicture}}
Left: plain green dots \code!g1!, \code!g2! (one input, two outputs) and pink dots \code!p1!,
\code!p2! (two inputs, one output), output $i$ of \code!g!$_j$ into input $j$ of \code!p!$_i$
($i,j\in\{1,2\}$, counted from $1$); terminals $a,b$ at the green inputs and $c,e$ at the pink
outputs. Right: a pink dot with two inputs into a plain green dot with two outputs.
$\sigma=\nu_d(3)/\nu_d(3)^2=d^{-1/2}$. Fixed sides; generated converse \cmplr{zxw_b2_converse}.

\cmplhead{zxw_k0}{[P1] K0: a pink $K_j$ state is copied by a plain green dot.}
\cmpleq{\begin{tikzpicture}[zxd]
\node[cmplpk] (p) at (0,0) {$K_j$}; \node[zdot] (g) at (0,0.8) {}; \draw[wire] (p) -- (g);
\draw[wire] (g) -- (-0.4,1.5); \draw[wire] (g) -- (0.4,1.5);
\end{tikzpicture}}{}{\begin{tikzpicture}[zxd]
\node[cmplpk] (p1) at (-0.45,0.4) {$K_j$}; \node[cmplpk] (p2) at (0.45,0.4) {$K_j$};
\draw[wire] (p1) -- (-0.45,1.5); \draw[wire] (p2) -- (0.45,1.5);
\end{tikzpicture}}
Left: a pink state with \py!Kj("j")! into a plain green dot with two outputs; right: two pink
states with \py!Kj("j")!. $\sigma=\nu_d(1)^2/\nu_d(1)=d^{-1/2}$. Fixed sides; promoted;
quantifiers \code!j!, \code!d!; generated converse \cmplr{zxw_k0_converse}.

\cmplhead{zxw_k1}{[P1] K1: a pink $K_j$ passes through a plain green dot.}
\cmpleq{\begin{tikzpicture}[zxd]
\node[cmplpk] (p) at (0,0) {$K_j$}; \node[zdot] (g) at (0,0.8) {};
\draw[wire] (0,-0.6) -- (p) -- (g);
\draw[wire] (g) -- (-0.5,1.5); \draw[wire] (g) -- (0.5,1.5); \node[dimlabel] at (0,1.4) {$\cdots n$};
\end{tikzpicture}}{}{\begin{tikzpicture}[zxd]
\node[zdot] (g) at (0,0) {};
\node[cmplpk] (q1) at (-0.5,0.85) {$K_j$}; \node[cmplpk] (q2) at (0.5,0.85) {$K_j$};
\draw[wire] (0,-0.6) -- (g); \draw[wire] (g) -- (q1) -- (-0.5,1.5); \draw[wire] (g) -- (q2) -- (0.5,1.5);
\node[dimlabel] at (0,0.85) {$\cdots$};
\end{tikzpicture}}
Left: a pink $K_j$ (one input, one output) into a plain green dot with one input and $n$ outputs.
Right: the plain green dot with a pink $K_j$ on each output. $\sigma=1$. Left arities: the anchor
must be an X spider with one input and one output; yield $n=(\text{port count})-1$ for each
distinct port count of an adjacent Z spider, in increasing order. Right arities: for a Z anchor,
$n=(\text{port count})-1$. Quantifiers \code!n!, \code!j!, \code!d!; generated converse
\cmplr{zxw_k1_converse}.

\cmplhead{zxw_k2}{[P1] K2, [P2] K2: a pink $K_j$ passes through a Z box.}
\cmpleq{\begin{tikzpicture}[zxd]
\node[cmplpk] (p) at (0,0) {$K_j$}; \node[cmplbox] (a) at (0,0.85) {$a$};
\draw[wire] (0,-0.6) -- (p) -- (a) -- (0,1.45);
\end{tikzpicture}}{a_{-j}\cdot}{\begin{tikzpicture}[zxd]
\node[cmplbox] (a) at (0,0) {$k_j(a)$}; \node[cmplpk] (p) at (0,0.85) {$K_j$};
\draw[wire] (0,-0.6) -- (a) -- (p) -- (0,1.45);
\end{tikzpicture}}
Left: a pink $K_j$ into a Z spider with \py!Var("a")!. Right: a Z spider with the derived
\py!shifted_vector(a, j)! into a pink $K_j$; the right template carries the scalar $a_{-j}$
(\py!_coefficient(a, -j)!, index read modulo $d$; $1$ for a phaseless $a$). $\sigma=a_{-j}$.
Condition \cmplr{a_minus_j_nonzero}: $a_{d-j}\neq 0$ by \py!_nonzero!. \py!concrete_only!;
promoted. The binding $j$ must be a concrete integer; otherwise \py!_int! raises
\py!RewriteDomainError! \code!"expected a concrete integer, got ..."!. Converse:
\cmplr{zxw_k2_converse}.

\cmplhead{zxw_zer}{[P1] Zer: the Z box $0$ state is the pink dot state.}
\cmpleq{\begin{tikzpicture}[zxd]
\node[cmplbox] (a) at (0,0) {$0$}; \draw[wire] (a) -- (0,0.8);
\end{tikzpicture}}{}{\begin{tikzpicture}[zxd]
\node[cmplpd] (p) at (0,0) {}; \draw[wire] (p) -- (0,0.8);
\end{tikzpicture}}
Left: a Z spider with no inputs, one output, and the fixed vector \py!ZERO!; right: a pink dot
state. $\sigma=\nu_d(1)=d^{-1/2}$. Fixed sides; promoted; generated converse
\cmplr{zxw_zer_converse}.

\cmplhead{zxw_p1}{[P1] P1 with (Du): $H$ followed by $H$ is $D$, the bent pink dot.}
\cmpleq{\begin{tikzpicture}[zxd]
\node[fbox] (f1) at (0,0) {$F$}; \node[fbox] (f2) at (0,0.7) {$F$};
\draw[wire] (0,-0.6) -- (f1) -- (f2) -- (0,1.3);
\end{tikzpicture}}{}{\begin{tikzpicture}[zxd]
\node[cmplpd] (p) at (0,0.55) {};
\draw[wire] (p) -- (0,1.3);
\draw[wire] (p) to[out=-20,in=90] (0.45,0.1) to[out=-90,in=0] (0.2,-0.25) to[out=180,in=90] (0,-0.6);
\end{tikzpicture}}
Left: two $F$ in series (terminals \code!in!, \code!out!). Right: a pink dot with no inputs and two
outputs, output $0$ the terminal \code!in! (bent down) and output $1$ the terminal \code!out!.
$\nu_d(2)=1$, so $\sigma=1$. Fixed sides; generated converse \cmplr{zxw_p1_converse}.

\cmplhead{zxw_p1_multiplier}{[P1] P1, [P2] P1: $H$ followed by $H$ is the multiplier $-1$.}
\cmpleq{\begin{tikzpicture}[zxd]
\node[fbox] (f1) at (0,0) {$F$}; \node[fbox] (f2) at (0,0.7) {$F$};
\draw[wire] (0,-0.6) -- (f1) -- (f2) -- (0,1.3);
\end{tikzpicture}}{}{\begin{tikzpicture}[zxd]
\node[cmplmult] (m) at (0,0.35) {$d-1$}; \draw[wire] (0,-0.6) -- (m) -- (0,1.3);
\end{tikzpicture}}
Left: two $F$ in series, defined at every dimension. Right: the multiplier labeled $d-1$ at a
concrete dimension (\py!None! at a symbolic one), with \py!min_size! $d-1$.
$\sigma=\nu_d(d)=d^{(d-2)/2}$. \py!concrete_only!; generated converse
\cmplr{zxw_p1_multiplier_converse}.

\cmplhead{zxw_h1}{[P1] H1: $H$ and $H^\dagger$ pass through a green circle spider $K_j$, negating $j$.}
\cmpleq{\begin{tikzpicture}[zxd]
\node[zs] (p) at (0,0.7) {$K_j$};
\node[fbox] (a1) at (-0.5,-0.2) {$F^3$}; \node[fbox] (a2) at (0.5,-0.2) {$F^3$};
\node[fbox] (b1) at (-0.5,1.6) {$F$}; \node[fbox] (b2) at (0.5,1.6) {$F$};
\draw[wire] (-0.5,-0.8) -- (a1) -- (p); \draw[wire] (0.5,-0.8) -- (a2) -- (p);
\draw[wire] (p) -- (b1) -- (-0.5,2.2); \draw[wire] (p) -- (b2) -- (0.5,2.2);
\node[dimlabel] at (0,-0.2) {$\cdots n$}; \node[dimlabel] at (0,1.6) {$\cdots m$};
\end{tikzpicture}}{}{\begin{tikzpicture}[zxd]
\node[zs] (p) at (0,0.7) {$K_{-j}$};
\node[fbox] (a1) at (-0.5,-0.2) {$F$}; \node[fbox] (a2) at (0.5,-0.2) {$F$};
\node[fbox] (b1) at (-0.5,1.6) {$F^3$}; \node[fbox] (b2) at (0.5,1.6) {$F^3$};
\draw[wire] (-0.5,-0.8) -- (a1) -- (p); \draw[wire] (0.5,-0.8) -- (a2) -- (p);
\draw[wire] (p) -- (b1) -- (-0.5,2.2); \draw[wire] (p) -- (b2) -- (0.5,2.2);
\node[dimlabel] at (0,-0.2) {$\cdots n$}; \node[dimlabel] at (0,1.6) {$\cdots m$};
\end{tikzpicture}}
A box $F^3$ is three $F$ in series. Left: a Z spider with $n$ inputs, $m$ outputs, and
\py!Kj("j")!, three $F$ before each input and one $F$ after each output. Right: the Z spider with
\py!Kj("j", sign=-1)!, one $F$ before each input and three after each output. $\sigma=1$. Both
sides: for a Z anchor with $p$ ports, yield $(n,m)=(0,p),(1,p-1),\dots,(p,0)$. Quantifiers
\code!n!, \code!m!, \code!j!, \code!d!; generated converse \cmplr{zxw_h1_converse}.

\cmplhead{zxw_pcy}{[P1] Pcy: a Z box on the input of a W node is copied onto both outputs.}
\cmpleq{\begin{tikzpicture}[zxd]
\node[cmplbox] (a) at (0,0) {$a$}; \node[wnode] (w) at (0,0.75) {};
\draw[wire] (0,-0.6) -- (a) -- (w); \draw[wire] (w) -- (-0.4,1.45); \draw[wire] (w) -- (0.4,1.45);
\end{tikzpicture}}{}{\begin{tikzpicture}[zxd]
\node[wnode] (w) at (0,0) {}; \node[cmplbox] (a1) at (-0.45,0.8) {$a$}; \node[cmplbox] (a2) at (0.45,0.8) {$a$};
\draw[wire] (0,-0.6) -- (w); \draw[wire] (w) -- (a1) -- (-0.45,1.45); \draw[wire] (w) -- (a2) -- (0.45,1.45);
\end{tikzpicture}}
Left: a Z spider (one input, one output, \py!Var("a")!) into a W node with two outputs. Right: the
W node with a Z spider \py!Var("a")! on each output. $\sigma=1$. Fixed sides; generated converse
\cmplr{zxw_pcy_converse}, whose matched side binds $a$ twice, so both boxes must carry the same
vector.

\cmplhead{zxw_ad}{[P1] AD, [P2] AD: Z box effects on a W node add.}
\cmpleq{\begin{tikzpicture}[zxd]
\node[wnode] (w) at (0,0) {}; \node[cmplbox] (a) at (-0.45,0.8) {$a$}; \node[cmplbox] (b) at (0.45,0.8) {$b$};
\draw[wire] (0,-0.6) -- (w); \draw[wire] (w) -- (a); \draw[wire] (w) -- (b);
\end{tikzpicture}}{}{\begin{tikzpicture}[zxd]
\node[cmplbox] (c) at (0,0.4) {$a+b$}; \draw[wire] (0,-0.6) -- (c);
\end{tikzpicture}}
Left: a W node with two outputs into Z effects \py!Var("a")! and \py!Var("b")!. Right: one Z
effect with the derived \py!sum_vector(a, b)!. $\sigma=1$. Fixed sides; the rule also applies at
a symbolic dimension when both vectors are rule-only; promoted. Converse: \cmplr{zxw_ad_converse}.

\cmplhead{zxw_sym}{[P1] Sym: the two outputs of a W node may be swapped.}
\cmpleq{\begin{tikzpicture}[zxd]
\node[wnode] (w) at (0,0) {}; \draw[wire] (0,-0.6) -- (w);
\draw[wire] (w) -- (-0.4,0.7) node[above,dimlabel] {$a$}; \draw[wire] (w) -- (0.4,0.7) node[above,dimlabel] {$b$};
\end{tikzpicture}}{}{\begin{tikzpicture}[zxd]
\node[wnode] (w) at (0,0) {}; \draw[wire] (0,-0.6) -- (w);
\draw[wire] (w.north) to[out=60,in=-90] (0.4,0.7) node[above,dimlabel] {$a$};
\draw[wire] (w.north) to[out=120,in=-90] (-0.4,0.7) node[above,dimlabel] {$b$};
\end{tikzpicture}}
Left: a W node whose outputs $0$ and $1$ are the terminals $a$ and $b$; right: the same node with
the terminals exchanged. $\sigma=1$. Fixed sides; generated converse \cmplr{zxw_sym_converse}.
The matcher identifies the outputs of a W node up to permutation, so the rule changes only the
port order of the result.

\cmplhead{zxw_aso}{[P1] Aso: the W node is associative.}
\cmpleq{\begin{tikzpicture}[zxd]
\node[wnode] (w1) at (0,0) {}; \node[wnode] (w2) at (-0.4,0.8) {};
\draw[wire] (0,-0.6) -- (w1); \draw[wire] (w1) -- (w2); \draw[wire] (w1) -- (0.5,1.6) node[above,dimlabel] {$c$};
\draw[wire] (w2) -- (-0.8,1.6) node[above,dimlabel] {$a$}; \draw[wire] (w2) -- (-0.1,1.6) node[above,dimlabel] {$b$};
\end{tikzpicture}}{}{\begin{tikzpicture}[zxd]
\node[wnode] (w1) at (0,0) {}; \node[wnode] (w2) at (0.4,0.8) {};
\draw[wire] (0,-0.6) -- (w1); \draw[wire] (w1) -- (w2); \draw[wire] (w1) -- (-0.5,1.6) node[above,dimlabel] {$a$};
\draw[wire] (w2) -- (0.1,1.6) node[above,dimlabel] {$b$}; \draw[wire] (w2) -- (0.8,1.6) node[above,dimlabel] {$c$};
\end{tikzpicture}}
Left: W node \code!w1! (two outputs), output $0$ into W node \code!w2! (two outputs); terminals
$a,b$ at the outputs of \code!w2! and $c$ at output $1$ of \code!w1!. Right: output $1$ of
\code!w1! into \code!w2!; $a$ at output $0$ of \code!w1!, $b,c$ at the outputs of \code!w2!.
$\sigma=1$. Fixed sides; promoted; generated converse \cmplr{zxw_aso_converse}.

\cmplhead{zxw_bzw}{[P1] BZW: the bialgebra of the W node and the plain green dot.}
\cmpleq{\begin{tikzpicture}[zxd]
\node[wnode] (w1) at (-0.5,0) {}; \node[wnode] (w2) at (0.5,0) {};
\node[zdot] (g1) at (-0.5,1.1) {}; \node[zdot] (g2) at (0.5,1.1) {};
\draw[wire] (w1) -- (g1); \draw[wire] (w1) -- (g2); \draw[wire] (w2) -- (g1); \draw[wire] (w2) -- (g2);
\draw[wire] (w1) -- (-0.5,-0.6) node[below,dimlabel] {$a$}; \draw[wire] (w2) -- (0.5,-0.6) node[below,dimlabel] {$b$};
\draw[wire] (g1) -- (-0.5,1.7) node[above,dimlabel] {$c$}; \draw[wire] (g2) -- (0.5,1.7) node[above,dimlabel] {$e$};
\end{tikzpicture}}{}{\begin{tikzpicture}[zxd]
\node[zdot] (g) at (0,0.2) {}; \node[wnode] (w) at (0,0.9) {}; \draw[wire] (g) -- (w);
\draw[wire] (g) -- (-0.4,-0.6) node[below,dimlabel] {$a$}; \draw[wire] (g) -- (0.4,-0.6) node[below,dimlabel] {$b$};
\draw[wire] (w) -- (-0.4,1.7) node[above,dimlabel] {$c$}; \draw[wire] (w) -- (0.4,1.7) node[above,dimlabel] {$e$};
\end{tikzpicture}}
Left: W nodes \code!w1!, \code!w2! (two outputs each) and plain green dots \code!g1!, \code!g2!
(two inputs, one output), output $i$ of \code!w!$_j$ into input $j$ of \code!g!$_i$. Right: a
plain green dot with two inputs into a W node with two outputs. $\sigma=1$. Fixed sides;
promoted; generated converse \cmplr{zxw_bzw_converse}.

\cmplhead{zxw_ww}{[P1] WW: the crossed bialgebra of W nodes and W transposes.}
\cmpleq{\begin{tikzpicture}[zxd]
\node[cmplwt] (t) at (-1,-0.35) {}; \node[wnode] (w) at (-1,0.35) {};
\node[wnode] (p) at (0.6,0.35) {};
\node[zdot] (g1) at (-0.6,1.35) {}; \node[zdot] (g2) at (0.4,1.35) {};
\draw[wire] (-1.35,-1.1) node[below,dimlabel] {$x_0$} -- (t); \draw[wire] (-0.65,-1.1) node[below,dimlabel] {$x_1$} -- (t);
\draw[wire] (t) -- (w); \draw[wire] (0.6,-1.1) node[below,dimlabel] {$u$} -- (p);
\draw[wire] (w) -- (g1); \draw[wire] (w) -- (g2); \draw[wire] (p) -- (g1); \draw[wire] (p) -- (g2);
\draw[wire] (g1) -- (-0.6,2) node[above,dimlabel] {$y_0$}; \draw[wire] (g2) -- (0.4,2) node[above,dimlabel] {$y_1$};
\end{tikzpicture}}{}{\begin{tikzpicture}[zxd]
\node[wnode] (w0) at (-1.5,-0.45) {}; \node[wnode] (w1) at (-0.5,-0.45) {};
\node[cmplwt] (t0) at (-1.5,0.55) {}; \node[cmplwt] (t1) at (-0.5,0.55) {};
\node[wnode] (p) at (0.7,0.0) {};
\node[zdot] (g0) at (-1.0,1.45) {}; \node[zdot] (g1) at (0.2,1.45) {};
\draw[wire] (-1.5,-1.1) node[below,dimlabel] {$x_0$} -- (w0); \draw[wire] (-0.5,-1.1) node[below,dimlabel] {$x_1$} -- (w1);
\draw[wire] (0.7,-1.1) node[below,dimlabel] {$u$} -- (p);
\draw[wire] (w0) -- (t0); \draw[wire] (w0) -- (t1); \draw[wire] (w1) -- (t0); \draw[wire] (w1) -- (t1);
\draw[wire] (p) -- (g0); \draw[wire] (p) -- (g1); \draw[wire] (t0) -- (g0); \draw[wire] (t1) -- (g1);
\draw[wire] (g0) -- (-1.0,2.05) node[above,dimlabel] {$y_0$}; \draw[wire] (g1) -- (0.2,2.05) node[above,dimlabel] {$y_1$};
\end{tikzpicture}}
Shown for $n=m=2$. Left: a W node \code!t! with $n$ outputs whose input is wired to the input of
a W node \code!w! with $m$ outputs, so \code!t! is a W transpose with the $n$ inputs $x_k$; a W
node \code!p! with $m$ outputs and input $u$; plain green dots \code!g!$_k$ (two inputs, output
$y_k$) with input $0$ from output $k$ of \code!p! and input $1$ from output $k$ of \code!w!.
Right: \code!p! with input $u$; for each $k<n$ a W node \code!w!$_k$ with $m$ outputs and input
$x_k$; for each $c<m$ a W node \code!t!$_c$ with $n$ outputs, output $k$ wired to output $c$ of
\code!w!$_k$ (a W transpose), the input of \code!t!$_c$ into input $1$ of \code!g!$_c$, and output
$c$ of \code!p! into input $0$ of \code!g!$_c$. $\sigma=1$. Left arities: when the anchor is a W
node whose input is wired to the input of another W node, yield $n=$ the anchor's output count
and $m=$ the other's. Right arities: for a W anchor with $m$ outputs, yield
$n=1,\dots,|\text{nodes}|$. Quantifiers \code!n!, \code!m!, \code!d!; generated converse
\cmplr{zxw_ww_converse}.

\cmplhead{zxw_bs0}{[P1] Bs0: the pink dot state into a W node is two pink dot states.}
\cmpleq{\begin{tikzpicture}[zxd]
\node[cmplpd] (p) at (0,0) {}; \node[wnode] (w) at (0,0.7) {}; \draw[wire] (p) -- (w);
\draw[wire] (w) -- (-0.4,1.4); \draw[wire] (w) -- (0.4,1.4);
\end{tikzpicture}}{}{\begin{tikzpicture}[zxd]
\node[cmplpd] (p1) at (-0.4,0.4) {}; \node[cmplpd] (p2) at (0.4,0.4) {};
\draw[wire] (p1) -- (-0.4,1.4); \draw[wire] (p2) -- (0.4,1.4);
\end{tikzpicture}}
$\sigma=\nu_d(1)^2/\nu_d(1)=d^{-1/2}$. Fixed sides; promoted; generated converse
\cmplr{zxw_bs0_converse}.

\cmplhead{zxw_bsj}{[P1] Bsj, [P2] Bsj: a pink $K_j$ state into a W node with one output on a plain green effect.}
\cmpleq{\begin{tikzpicture}[zxd]
\node[cmplpk] (p) at (0,0) {$K_j$}; \node[wnode] (w) at (0,0.75) {}; \node[zdot] (e) at (-0.45,1.35) {};
\draw[wire] (p) -- (w); \draw[wire] (w) -- (e); \draw[wire] (w) -- (0.45,1.6);
\end{tikzpicture}}{}{\begin{tikzpicture}[zxd]
\node[cmplbox] (z) at (0,0.4) {$u_{-j}$}; \draw[wire] (z) -- (0,1.4);
\end{tikzpicture}}
Left: a pink $K_j$ state into a W node with two outputs, output $0$ into a plain green effect,
output $1$ the terminal. Right: a Z state with the derived \py!unit_vector(-j)!, the state
$T_j$ of the papers: $u_{-j}$, with coefficient $1$ at $0$ and at $-j\bmod d$ and $0$ elsewhere.
$\sigma=\nu_d(1)^{-1}=\sqrt d$. Fixed sides; promoted; quantifiers \code!j!, \code!d!. Converse:
\cmplr{zxw_bsj_converse}.

\cmplhead{zxw_ta}{[P1] TA: the trialgebra rule; the pink dot becomes a W node.}
\cmpleq{\begin{tikzpicture}[zxd]
\node[wnode] (w) at (-0.5,0) {}; \node[cmplpd] (p) at (0.5,0) {};
\node[zdot] (g1) at (-0.5,1.1) {}; \node[zdot] (g2) at (0.5,1.1) {};
\draw[wire] (w) -- (g1); \draw[wire] (w) -- (g2); \draw[wire] (p) -- (g1); \draw[wire] (p) -- (g2);
\draw[wire] (w) -- (-0.5,-0.6) node[below,dimlabel] {$a$}; \draw[wire] (p) -- (0.5,-0.6) node[below,dimlabel] {$b$};
\draw[wire] (g1) -- (-0.5,1.7) node[above,dimlabel] {$c$}; \draw[wire] (g2) -- (0.5,1.7) node[above,dimlabel] {$e$};
\end{tikzpicture}}{}{\begin{tikzpicture}[zxd]
\node[wnode] (w) at (-0.5,0) {}; \node[wnode] (p) at (0.5,0) {};
\node[zdot] (g1) at (-0.5,1.1) {}; \node[zdot] (g2) at (0.5,1.1) {};
\draw[wire] (w) -- (g1); \draw[wire] (w) -- (g2); \draw[wire] (p) -- (g1); \draw[wire] (p) -- (g2);
\draw[wire] (w) -- (-0.5,-0.6) node[below,dimlabel] {$a$}; \draw[wire] (p) -- (0.5,-0.6) node[below,dimlabel] {$b$};
\draw[wire] (g1) -- (-0.5,1.7) node[above,dimlabel] {$c$}; \draw[wire] (g2) -- (0.5,1.7) node[above,dimlabel] {$e$};
\end{tikzpicture}}
Left (\py!_ta_left(False)!): a W node \code!w! and a pink dot \code!p! (one input, two outputs);
output $i$ of \code!w! into input $0$ of the plain green dot \code!g!$_{i+1}$ and output $i$ of
\code!p! into its input $1$. Right (\py!_ta_left(True)!): the same with \code!p! a W node.
$\sigma=\nu_d(3)^{-1}=d^{-1/2}$. Fixed sides; promoted; generated converse \cmplr{zxw_ta_converse}.

\cmplhead{zxw_hd}{[P1] HD: the decomposition of the Hadamard box.}
\cmpleq{\begin{tikzpicture}[zxd]
\node[cmplpd] (top) at (0,-2.7) {};
\node[cmplpk] (k1) at (0,-2.0) {$K_1$}; \node[wnode] (w) at (0,-1.3) {};
\node[cmplmultt] (m0) at (-1.3,-1.3) {$0$}; \node[cmplmultt] (m1) at (1.3,-1.3) {$1$};
\node[zdot] (g0) at (-1.3,-0.35) {}; \node[zdot] (g1) at (1.3,-0.35) {};
\node[cmplv] (M0) at (-1.3,0.45) {$M$}; \node[cmplv] (M1) at (1.3,0.45) {$M$};
\node[zs] (K0) at (-1.3,1.3) {$K_0$}; \node[zs] (Kd) at (1.3,1.3) {$K_{d-1}$};
\node[cmplpd] (bot) at (0,2.1) {};
\draw[wire] (0,-3.2) -- (top); \draw[wire] (top) -- (m0); \draw[wire] (top) -- (m1);
\draw[wire] (k1) -- (w); \draw[wire] (w) -- (g0); \draw[wire] (w) -- (g1);
\draw[wire] (m0) -- (g0) -- (M0) -- (K0) -- (bot); \draw[wire] (m1) -- (g1) -- (M1) -- (Kd) -- (bot);
\draw[wire] (bot) -- (0,2.6);
\node[dimlabel] at (0,0.45) {$\cdots$}; \node[dimlabel] at (0,1.3) {$\cdots$};
\end{tikzpicture}}{}{\begin{tikzpicture}[zxd]
\node[fbox] (f) at (0,0) {$F$}; \draw[wire] (0,-0.6) -- (f) -- (0,0.6);
\end{tikzpicture}}
At concrete $d$ the left side is $d^{-1/2}$ (its template scalar) times the following diagram,
drawn with the columns $c=0$ and $c=d-1$. A pink dot \code!top! takes the rule's input and has
$d$ outputs. For each column $c=0,\dots,d-1$: output $c$ of \code!top! passes through the
transposed multiplier labeled $(-c)\bmod d$ into input $0$ of a plain green dot \code!c!$_c$ with
three inputs and no outputs; input $1$ of \code!c!$_c$ comes from output $c$ of a W node with $d$
outputs fed by a pink $K_1$ state; input $2$ of \code!c!$_c$ is wired to the input of an M box, so
\code!c!$_c$ acts as a dot with two inputs and one output into the M box; the M box feeds a Z
spider with phase vector $K_c$ (phaseless for $c=0$), which feeds input $c$ of a pink dot
\code!bottom! with $d$ inputs whose output is the rule's output. The right side is one $F$. With
$E_L=(d-2)(d^2+2d+3)/4$ the exponent of $N_L$,
\[
\sigma \;=\; d^{\,1/2-E_L}\qquad(\sigma=\sqrt2,\ 3^{-4},\ 4^{-13},\ 5^{-28}\ \text{at}\ d=2,3,4,5).
\]
\py!concrete_only!; the left side grows; promoted; generated converse \cmplr{zxw_hd_converse}.

\cmplhead{zxw_zv}{[P1] ZV: a Z box state into V is a constant Z box state.}
\cmpleq{\begin{tikzpicture}[zxd]
\node[cmplbox] (a) at (0,0) {$a$}; \node[cmplv] (v) at (0,0.8) {$V$};
\draw[wire] (a) -- (v) -- (0,1.5);
\end{tikzpicture}}{}{\begin{tikzpicture}[zxd]
\node[cmplbox] (c) at (0,0.4) {$(a_{d-1},\dots,a_{d-1})$}; \draw[wire] (c) -- (0,1.5);
\end{tikzpicture}}
Left: a Z state with \py!Var("a")! into the V box. Right: a Z state with the derived
\py!constant_vector(a_{-1})!. With the V-box normalization,
$\sigma=d^{-(d/2-1)-(d-1)(d-2)/4}$ ($1$, $1/3$, $1/32$, $\sqrt5/3125$ at $d=2,3,4,5$).
\py!concrete_only!; the left side grows; promoted. Converse: \cmplr{zxw_zv_converse}.

\cmplhead{zxw_va}{[P1] VA: boxes on a W node through multipliers into a pink dot are V followed by a Z box.}
\cmpleq{\begin{tikzpicture}[zxd]
\node[wnode] (w) at (0,-0.9) {};
\node[cmplbox] (b1) at (-0.9,-0.1) {$a_1$}; \node[cmplbox] (b2) at (0.9,-0.1) {$a_{d-1}$};
\node[cmplmult] (m1) at (-0.9,0.7) {$-1$}; \node[cmplmult] (m2) at (0.9,0.7) {$1$};
\node[cmplpd] (q) at (0,1.5) {};
\draw[wire] (0,-1.5) -- (w); \draw[wire] (w) -- (b1) -- (m1) -- (q); \draw[wire] (w) -- (b2) -- (m2) -- (q);
\draw[wire] (q) -- (0,2.0); \node[dimlabel] at (0,0.3) {$\cdots$};
\end{tikzpicture}}{}{\begin{tikzpicture}[zxd]
\node[cmplv] (v) at (0,0) {$V$}; \node[cmplbox] (a) at (0,0.8) {$a$};
\draw[wire] (0,-0.6) -- (v) -- (a) -- (0,1.4);
\end{tikzpicture}}
In this picture a box labeled $a_k$ is the box labeled $a_k$, $(0,\dots,0,a_k)$. Left: a W
node with $d-1$ outputs; output $k-1$ into a Z spider with \py!LastEntry("a_k")!, then the
multiplier labeled $(-k)\bmod d$, into input $k-1$ of a pink dot with $d-1$ inputs. Right: the V
box into a Z spider with the derived \py!vector_of([a_1, ..., a_{d-1}])!, the Z box
$a=(a_1,\dots,a_{d-1})$. $\sigma=1$, since both sides carry the V-box normalization.
\py!concrete_only!; both sides grow; promoted. Converse: \cmplr{zxw_va_converse}.

\cmplhead{zxw_vw}{[P1] VW: V into a W node is the W node with V on both outputs.}
\cmpleq{\begin{tikzpicture}[zxd]
\node[cmplv] (v) at (0,0) {$V$}; \node[wnode] (w) at (0,0.75) {};
\draw[wire] (0,-0.6) -- (v) -- (w); \draw[wire] (w) -- (-0.4,1.45); \draw[wire] (w) -- (0.4,1.45);
\end{tikzpicture}}{}{\begin{tikzpicture}[zxd]
\node[wnode] (w) at (0,0) {}; \node[cmplv] (v1) at (-0.5,0.8) {$V$}; \node[cmplv] (v2) at (0.5,0.8) {$V$};
\draw[wire] (0,-0.6) -- (w); \draw[wire] (w) -- (v1) -- (-0.5,1.45); \draw[wire] (w) -- (v2) -- (0.5,1.45);
\end{tikzpicture}}
$\sigma=d^{\,d/2-1+(d-1)(d-2)/4}$ ($1$, $3$, $32$, $625\sqrt5$ at $d=2,3,4,5$). \py!concrete_only!;
both sides grow; generated converse \cmplr{zxw_vw_converse}.

\cmplhead{zxw_kz}{[P1] KZ: a pink $K_1$ state on a W node whose other outputs meet a pink dot.}
\cmpleq{\begin{tikzpicture}[zxd]
\node[cmplpk] (p) at (-0.6,-1.1) {$K_1$}; \node[wnode] (w) at (-0.6,-0.4) {};
\node[cmplpd] (q) at (1.0,-1.1) {};
\node[cmplmultt] (e1) at (0.5,-0.2) {$e_1$}; \node[cmplmultt] (e2) at (1.5,-0.2) {$e_2$};
\node[zdot] (g1) at (0.1,0.7) {}; \node[zdot] (g2) at (1.1,0.7) {};
\draw[wire] (p) -- (w); \draw[wire] (w) -- (-1.1,1.3) node[above,dimlabel] {$o_0$};
\draw[wire] (w) -- (g1); \draw[wire] (w) -- (g2);
\draw[wire] (q) -- (e1) -- (g1); \draw[wire] (q) -- (e2) -- (g2);
\draw[wire] (g1) -- (0.1,1.3) node[above,dimlabel] {$o_1$}; \draw[wire] (g2) -- (1.1,1.3) node[above,dimlabel] {$o_2$};
\end{tikzpicture}}{}{\begin{tikzpicture}[zxd]
\node[cmplpk] (p) at (-0.7,0.3) {$K_1$}; \node[cmplpd] (q1) at (0.1,0.3) {}; \node[cmplpd] (q2) at (0.8,0.3) {};
\draw[wire] (p) -- (-0.7,1.3) node[above,dimlabel] {$o_0$}; \draw[wire] (q1) -- (0.1,1.3) node[above,dimlabel] {$o_1$};
\draw[wire] (q2) -- (0.8,1.3) node[above,dimlabel] {$o_2$};
\end{tikzpicture}}
Shown for $n=2$. Left: a W node with $n+1$ outputs fed by a pink $K_1$ state (\py!k_fixed(1)!);
output $0$ the terminal $o_0$; output $k+1$ into input $0$ of a plain green dot \code!g!$_k$ (two
inputs, output $o_{k+1}$); a pink dot with $n$ outputs, output $k$ through the transposed
multiplier labeled $e_k$ into input $1$ of \code!g!$_k$. Right: a pink $K_1$ state and $n$ pink
dot states. $\sigma=d^{\,1-n-\sum_k (e_k-1)/2}$. Left arities (dimension-aware): the anchor must
be a W node with $n+1\ge2$ outputs; collect the input counts of the Z spiders adjacent to the Z
neighbors of the anchor, keep those $e$ with $1\le e\le d-1$ (with no upper bound at a symbolic
dimension), and yield every multiset of $n$ of them (\py!combinations_with_replacement! of the
sorted labels) as $e_0\le\dots\le e_{n-1}$. Right arities: at a concrete $d$, for an X-spider
anchor with no inputs, let $s$ be the number of phaseless X spiders with no inputs and one
output; yield $n=1,\dots,s$ with every multiset of $n$ labels from $1,\dots,d-1$.
\py!concrete_only!; promoted; quantifiers \code!n!, \code!e_1..e_n!, \code!d!; generated
converse \cmplr{zxw_kz_converse}.


\section{The Axioms of [P2] Not in [P1]}\label{sec:completeness-p2}

The axioms S2, B2, K2, P1, AD, and Bsj of [P2] coincide with rules of \cref{sec:completeness-p1}.
This section states the remaining 16 rules: eight at one dimension (WF, (2.10), (2.9), MW, D1, WW,
VA, HD) and eight on wires of several dimensions ((2.8), DD, E, DZX, K0, BZW, TA, S1).

\subsection{Single-dimensional rules}

\cmplhead{qf_wf}{[P2] WF: a W node fed by an output of another is one W node.}
\cmpleq{\begin{tikzpicture}[zxd]
\node[wnode] (w1) at (0,0) {}; \node[wnode] (w2) at (0,0.8) {};
\draw[wire] (0,-0.6) -- (w1); \draw[wire] (w1) -- (w2);
\draw[wire] (w1) -- (-0.7,1.6) node[above,dimlabel] {$p$}; \draw[wire] (w1) -- (0.7,1.6) node[above,dimlabel] {$q$};
\draw[wire] (w2) -- (-0.25,1.6); \draw[wire] (w2) -- (0.25,1.6); \node[dimlabel] at (0,1.75) {$r$};
\end{tikzpicture}}{}{\begin{tikzpicture}[zxd]
\node[wnode] (w) at (0,0.3) {}; \draw[wire] (0,-0.6) -- (w);
\draw[wire] (w) -- (-0.7,1.6); \draw[wire] (w) -- (-0.25,1.6); \draw[wire] (w) -- (0.25,1.6); \draw[wire] (w) -- (0.7,1.6);
\node[dimlabel] at (0,1.75) {$p+r+q$};
\end{tikzpicture}}
Left: a W node \code!w1! with $p+1+q$ outputs whose output $p$ feeds a W node \code!w2! with $r$
outputs; the terminals $o_0,\dots,o_{p+r+q-1}$ are, in order, the first $p$ outputs of \code!w1!,
the $r$ outputs of \code!w2!, and the last $q$ outputs of \code!w1!. Right: one W node with
$p+q+r$ outputs, output $k$ the terminal $o_k$. $\sigma=1$. Left arities: for a W anchor with at
least one output, for each output wired to the input of a different W node, collect that node's
output count $r$; yield $(p,q,r)=(0,\text{outputs}-1,r)$ for each distinct $r$ in increasing
order. Right arities: for a W anchor with $t$ outputs, yield $(0,t-r,r)$ for $r=0,\dots,t$.
Quantifiers \code!p!, \code!q!, \code!r!, \code!d!; generated converse \cmplr{qf_wf_converse}.

\cmplhead{qf_w_sym}{[P2] (2.10): adjacent outputs of a W node may be swapped.}
\cmpleq{\begin{tikzpicture}[zxd]
\node[wnode] (w) at (0,0) {}; \draw[wire] (0,-0.6) -- (w);
\draw[wire] (w) -- (-0.7,0.9); \draw[wire] (w) -- (-0.2,0.9) node[above,dimlabel] {$k$};
\draw[wire] (w) -- (0.2,0.9) node[above,dimlabel] {$k{+}1$}; \draw[wire] (w) -- (0.7,0.9);
\end{tikzpicture}}{}{\begin{tikzpicture}[zxd]
\node[wnode] (w) at (0,0) {}; \draw[wire] (0,-0.6) -- (w);
\draw[wire] (w) -- (-0.7,0.9); \draw[wire] (w.north) to[out=70,in=-90] (0.2,0.9) node[above,dimlabel] {$k$};
\draw[wire] (w.north) to[out=110,in=-90] (-0.2,0.9) node[above,dimlabel] {$k{+}1$}; \draw[wire] (w) -- (0.7,0.9);
\end{tikzpicture}}
Left: a W node with $n$ outputs, terminal $o_p$ at output $p$. Right: the same with outputs $k$
and $k+1$ exchanged among the terminals. $\sigma=1$. Both sides: for a W anchor with $n$ outputs,
yield $(n,k)$ for $k=0,\dots,n-2$. Quantifiers \code!n!, \code!k!, \code!d!; generated converse
\cmplr{qf_w_sym_converse}.

\cmplhead{qf_z_flex}{[P2] (2.9): an input port of a Z box bends to an output port.}
\cmpleq{\begin{tikzpicture}[zxd]
\node[cmplbox] (z) at (0,0.4) {$a$};
\draw[wire] (z) -- (-0.4,-0.5); \draw[wire] (z) -- (0.4,-0.5) node[below,dimlabel] {$i_{m-1}$};
\draw[wire] (z) -- (-0.4,1.3); \draw[wire] (z) -- (0.4,1.3);
\end{tikzpicture}}{}{\begin{tikzpicture}[zxd]
\node[cmplbox] (z) at (0,0.4) {$a$};
\draw[wire] (z) -- (-0.4,-0.5); \draw[wire] (z) -- (-0.4,1.3); \draw[wire] (z) -- (0.2,1.3);
\draw[wire] (z) to[out=30,in=90] (0.75,0.8) -- (0.75,-0.5) node[below,dimlabel] {$i_{m-1}$};
\end{tikzpicture}}
Left: a Z spider with $m$ inputs, $n$ outputs, and \py!Var("a")!. Right: a Z spider with $m-1$
inputs and $n+1$ outputs and the same parameter, whose last output $n$ is the terminal
$i_{m-1}$. $\sigma=1$. Left arities: for a Z anchor with $m\ge1$ inputs and $n$ outputs, yield
$(m,n)$. Right arities: for a Z anchor with $n'\ge1$ outputs and $m'$ inputs, yield
$(m'+1,n'-1)$. Quantifiers \code!m!, \code!n!, \code!d!; generated converse
\cmplr{qf_z_flex_converse}.

\cmplhead{qf_mw}{[P2] MW: a multiplier through the box labeled 1 and a W node.}
\cmpleq{\begin{tikzpicture}[zxd]
\node[cmplbox] (b) at (0,0) {$1$}; \node[cmplmult] (m) at (0,0.75) {$e$}; \node[wnode] (w) at (0,1.45) {};
\draw[wire] (0,-0.6) -- (b) -- (m) -- (w); \draw[wire] (w) -- (-0.4,2.1); \draw[wire] (w) -- (0.4,2.1);
\end{tikzpicture}}{}{\begin{tikzpicture}[zxd]
\node[cmplbox] (b) at (0,0) {$1$}; \node[wnode] (w) at (0,0.7) {};
\node[cmplmult] (m1) at (-0.45,1.45) {$e$}; \node[cmplmult] (m2) at (0.45,1.45) {$e$};
\draw[wire] (0,-0.6) -- (b) -- (w); \draw[wire] (w) -- (m1) -- (-0.45,2.1); \draw[wire] (w) -- (m2) -- (0.45,2.1);
\end{tikzpicture}}
Left: the box labeled $1$ (\py!ONE_LAST!, one input, one output) into the multiplier labeled $e$
into a W node with two outputs. Right: the box into the W node with the multiplier labeled $e$ on
each output. $\sigma=\nu_d(e+1)^2/\nu_d(e+1)=d^{(e-1)/2}$. Left arities (dimension-aware): for a Z
anchor, collect the output counts of the adjacent Z spiders; yield each $e$ in increasing order
with $e\not\equiv 0\pmod d$ (every $e$ at a symbolic $d$). Right arities: for a Z anchor, for each
adjacent W node, collect the output counts of the Z spiders adjacent to that W node other than
the anchor; yield as on the left. Quantifiers \code!e!, \code!d!; generated converse
\cmplr{qf_mw_converse}.

\cmplhead{qf_d1}{[P2] D1: $D$ on every input of a Z box is $D$ on every output of the reversed box.}
\cmpleq{\begin{tikzpicture}[zxd]
\node[cmplbox] (z) at (0,0.7) {$a$};
\node[fbox] (f1) at (-0.45,-0.2) {$F^2$}; \node[fbox] (f2) at (0.45,-0.2) {$F^2$};
\draw[wire] (-0.45,-0.8) -- (f1) -- (z); \draw[wire] (0.45,-0.8) -- (f2) -- (z);
\draw[wire] (z) -- (-0.4,1.5); \draw[wire] (z) -- (0.4,1.5);
\node[dimlabel] at (0,-0.75) {$n$}; \node[dimlabel] at (0,1.5) {$m$};
\end{tikzpicture}}{}{\begin{tikzpicture}[zxd]
\node[cmplbox] (z) at (0,0.1) {$\overleftarrow{a}$};
\node[fbox] (f1) at (-0.45,1.0) {$F^2$}; \node[fbox] (f2) at (0.45,1.0) {$F^2$};
\draw[wire] (-0.4,-0.8) -- (z); \draw[wire] (0.4,-0.8) -- (z);
\draw[wire] (z) -- (f1) -- (-0.45,1.6); \draw[wire] (z) -- (f2) -- (0.45,1.6);
\node[dimlabel] at (0,-0.75) {$n$}; \node[dimlabel] at (0,1.6) {$m$};
\end{tikzpicture}}
A box $F^2$ is two $F$ in series. Left: a Z spider with $n$ inputs, $m$ outputs, and
\py!Var("a")!, two $F$ before each input. Right: a Z spider with \py!reversed_vector(a)!, two $F$
after each output. $\sigma=1$. Both sides: for a Z anchor with $p$ ports, yield $(n,m)=(n,p-n)$
for $n=1,\dots,p$. \py!concrete_only!; quantifiers \code!n!, \code!m!, \code!d!. Converse:
\cmplr{qf_d1_converse}.

\cmplhead{qf_ww}{[P2] WW: a W transpose into a W node.}
\cmpleq{\begin{tikzpicture}[zxd]
\node[cmplwt] (t) at (0,0) {}; \node[wnode] (w) at (0,0.7) {};
\draw[wire] (-0.4,-0.7) node[below,dimlabel] {$x_0$} -- (t); \draw[wire] (0.4,-0.7) node[below,dimlabel] {$x_1$} -- (t);
\draw[wire] (t) -- (w); \draw[wire] (w) -- (-0.4,1.4) node[above,dimlabel] {$y_0$}; \draw[wire] (w) -- (0.4,1.4) node[above,dimlabel] {$y_1$};
\end{tikzpicture}}{}{\begin{tikzpicture}[zxd]
\node[wnode] (u0) at (-1.3,-0.5) {}; \node[wnode] (u1) at (-0.3,-0.5) {};
\node[cmplwt] (v0) at (-1.3,0.5) {}; \node[cmplwt] (v1) at (-0.3,0.5) {};
\node[zdot] (g0) at (-1.3,1.3) {}; \node[zdot] (g1) at (-0.3,1.3) {};
\node[cmplwt] (s) at (0.8,2.0) {}; \node[zdot] (e) at (0.8,2.7) {};
\draw[wire] (-1.3,-1.2) node[below,dimlabel] {$x_0$} -- (u0); \draw[wire] (-0.3,-1.2) node[below,dimlabel] {$x_1$} -- (u1);
\draw[wire] (u0) -- (v0); \draw[wire] (u0) -- (v1); \draw[wire] (u1) -- (v0); \draw[wire] (u1) -- (v1);
\draw[wire] (v0) -- (g0); \draw[wire] (v1) -- (g1);
\draw[wire] (g0) -- (-1.5,2.7) node[above,dimlabel] {$y_0$}; \draw[wire] (g1) -- (-0.5,2.7) node[above,dimlabel] {$y_1$};
\draw[wire] (g0) to[out=60,in=-120] (s); \draw[wire] (g1) to[out=60,in=-100] (s); \draw[wire] (s) -- (e);
\end{tikzpicture}}
Shown for $m=n=2$. Left: a W node \code!t! with $m$ outputs whose input is wired to the input of a
W node \code!w! with $n$ outputs; terminals $x_k$ at the outputs of \code!t! and $y_c$ at the
outputs of \code!w!. Right: for $k<m$, a W node \code!u!$_k$ with $n$ outputs and input $x_k$; a W
node \code!s! with $n$ outputs whose input is wired to the input of a plain green effect \code!e!;
for $c<n$, a W node \code!v!$_c$ with $m$ outputs, output $k$ wired to output $c$ of
\code!u!$_k$, the input of \code!v!$_c$ into a plain green dot \code!g!$_c$ with one input and two
outputs, output $0$ of \code!g!$_c$ the terminal $y_c$ and output $1$ wired to output $c$ of
\code!s!. So \code!v!$_c$ and \code!s! act as W transposes and \code!s! ends in the green effect.
$\sigma=1$. Left arities: when the anchor is a W node whose input is wired to the input of a
different W node, yield $m=$ the anchor's output count and $n=$ the other's. Right arities: for a W
anchor with $n$ outputs, for each output wired to an output of a W node, collect that node's output
count $m$; yield each distinct $(m,n)$ in increasing $m$. Quantifiers \code!m!, \code!n!, \code!d!;
generated converse \cmplr{qf_ww_converse}.

\cmplhead{qf_va}{[P2] VA: boxes on a W node through multipliers into a pink dot.}
\cmpleq{\begin{tikzpicture}[zxd]
\node[wnode] (w) at (0,-0.9) {};
\node[cmplbox] (b0) at (-1.3,-0.1) {$1$}; \node[cmplbox] (b1) at (0,-0.1) {$a_1$}; \node[cmplbox] (b2) at (1.3,-0.1) {$a_{d-1}$};
\node[cmplmult] (m0) at (-1.3,0.7) {$0$}; \node[cmplmult] (m1) at (0,0.7) {$-1$}; \node[cmplmult] (m2) at (1.3,0.7) {$1$};
\node[cmplpd] (q) at (0,1.5) {};
\draw[wire] (0,-1.5) -- (w); \draw[wire] (w) -- (b0) -- (m0) -- (q); \draw[wire] (w) -- (b1) -- (m1) -- (q);
\draw[wire] (w) -- (b2) -- (m2) -- (q);
\draw[wire] (q) -- (0,2.0); \node[dimlabel] at (0.65,0.3) {$\cdots$};
\end{tikzpicture}}{}{\begin{tikzpicture}[zxd]
\node[cmplbox] (b) at (0,-1.2) {$1$}; \node[cmplpk] (km) at (0,-0.5) {$K_{-1}$};
\node[cmplwt] (wt) at (0.3,0.25) {}; \node[zdot] (g) at (0.75,-0.5) {};
\node[cmplpk] (kp) at (0.3,0.95) {$K_1$}; \node[cmplbox] (a) at (0.3,1.65) {$a$};
\draw[wire] (0,-1.8) -- (b) -- (km); \draw[wire] (km) -- (wt); \draw[wire] (g) -- (wt);
\draw[wire] (wt) -- (kp) -- (a) -- (0.3,2.2);
\end{tikzpicture}}
Left: a W node with $d$ outputs; output $0$ into the box labeled $1$, then the multiplier labeled
$0$, into input $0$ of a pink dot with $d$ inputs; output $k\ge1$ into a Z spider with
\py!LastEntry("a_k")!, then the multiplier labeled $(-k)\bmod d$, into input $k$ of the pink dot.
Right: the box labeled $1$, a pink $K_{-1}$ (1 input, 1 output) whose output is wired to output
$0$ of a W node \code!wt! with two outputs, a plain green state wired to output $1$ of \code!wt!,
the input of \code!wt! wired to the input of a pink $K_1$, and the pink $K_1$ into a Z spider with
the derived \py!vector_of([a_1, ..., a_{d-1}])!. \code!wt! acts as a W transpose with two inputs.
$\sigma=d^{-(d-2)(d+1)/4}$ ($1$, $1/3$, $1/32$, $\sqrt5/3125$ at $d=2,3,4,5$). \py!concrete_only!;
the left side grows. Converse: \cmplr{qf_va_converse}.

\cmplhead{qf_hd}{[P2] HD: the decomposition of the Hadamard box.}
\cmpleq{\begin{tikzpicture}[zxd]
\node[cmplpd] (top) at (0,-2.7) {};
\node[cmplpk] (k1) at (0,-2.0) {$K_1$}; \node[wnode] (w) at (0,-1.3) {};
\node[cmplmultt] (m0) at (-1.5,-1.3) {$0$}; \node[cmplmultt] (m1) at (1.5,-1.3) {$1$};
\node[zdot] (g0) at (-1.5,-0.45) {}; \node[zdot] (g1) at (1.5,-0.45) {};
\node[cmplpk] (km0) at (-1.5,0.2) {$K_{-1}$}; \node[cmplpk] (km1) at (1.5,0.2) {$K_{-1}$};
\node[cmplwt] (wt0) at (-1.3,0.9) {}; \node[cmplwt] (wt1) at (1.7,0.9) {};
\node[zdot] (s0) at (-0.85,0.2) {}; \node[zdot] (s1) at (2.15,0.2) {};
\node[zs] (K0) at (-1.3,1.6) {$K_0$}; \node[zs] (Kd) at (1.7,1.6) {$K_{d-1}$};
\node[cmplpd] (bot) at (0.2,2.4) {};
\draw[wire] (0,-3.2) -- (top); \draw[wire] (top) -- (m0); \draw[wire] (top) -- (m1);
\draw[wire] (k1) -- (w); \draw[wire] (w) -- (g0); \draw[wire] (w) -- (g1);
\draw[wire] (m0) -- (g0) -- (km0) -- (wt0); \draw[wire] (s0) -- (wt0); \draw[wire] (wt0) -- (K0) -- (bot);
\draw[wire] (m1) -- (g1) -- (km1) -- (wt1); \draw[wire] (s1) -- (wt1); \draw[wire] (wt1) -- (Kd) -- (bot);
\draw[wire] (bot) -- (0.2,2.9);
\node[dimlabel] at (0.2,0.9) {$\cdots$};
\end{tikzpicture}}{}{\begin{tikzpicture}[zxd]
\node[fbox] (f) at (0,0) {$F$}; \draw[wire] (0,-0.6) -- (f) -- (0,0.6);
\end{tikzpicture}}
The left side is $(-1)^{d+1}/\sqrt d$ (its template scalar) times the diagram of \cmplr{zxw_hd}
with two changes per column $c$: the plain green dot \code!c!$_c$ has two inputs and one output,
and the M box is replaced by a pink $K_{-1}$ whose output is wired to output $1$ of a W node
\code!wt!$_c$ with two outputs, a plain green state wired to output $0$ of \code!wt!$_c$, and the
input of \code!wt!$_c$ wired to the input of the Z spider $K_c$. The right side is one $F$. With
$E_L=(d-2)(d+3)/4$,
\[
\sigma\;=\;(-1)^{d+1}\,d^{\,1/2-E_L}\qquad(-\sqrt2,\ 1/3,\ -1/64,\ \sqrt5/15625\ \text{at}\ d=2,3,4,5).
\]
\py!concrete_only!; the left side grows; generated converse \cmplr{qf_hd_converse}.

\subsection{Mixed-dimensional rules}\label{sec:completeness-mixed}

These rules read every port dimension from the template (\py!TNode.dim! or \py!port_dims!), so
their arities include dimensions. All are \py!concrete_only!. Their enumerators draw dimensions
from \py!_host_values!, so a rule instance whose dimensions do not all occur among the host's
concrete port dimensions is not tried, except for the anchor's own dimension. In the pictures a
port label is its dimension.

\cmplhead{binder_definition}{[P2] (2.8): the merger $B$ is the splitter $S$ with its ports bent.}
\cmpleq{\begin{tikzpicture}[zxd]
\node[bnode] (b) at (0,0.3) {$B$};
\draw[wire] (b.south west) -- ++(0,-0.6) node[below,dimlabel] {$s$}; \draw[wire] (b.south east) -- ++(0,-0.6) node[below,dimlabel] {$t$};
\draw[wire] (b) -- (0,1.2) node[above,dimlabel] {$st$};
\end{tikzpicture}}{}{\begin{tikzpicture}[zxd]
\node[snode] (s) at (0,0.3) {$S$};
\draw[wire] (s.north west) to[out=90,in=90] ($(s.north west)+(-0.55,0)$) -- ++(0,-1.0) node[below,dimlabel] {$s$};
\draw[wire] (s.north east) to[out=90,in=90] ($(s.north east)+(0.55,0)$) -- ++(0,-1.0) node[below,dimlabel] {$t$};
\draw[wire] (s.south) to[out=-90,in=-90] ($(s.south)+(0.0,-0.3)+(1.3,0)$) -- ++(0,1.6) node[above,dimlabel] {$st$};
\end{tikzpicture}}
Left: $B$ with input dimensions $(s,t)$ and output dimension $st$; terminals $a$, $b$ at its inputs
and $c$ at its output. Right: $S$ from $st$ to $(s,t)$ with $a$, $b$ at its outputs and $c$ at its
input. $\sigma=1$. Left arities: for a $B$ anchor whose two input dimensions resolve to concrete
values, yield $(s,t)$ from them. Right arities: for an $S$ anchor, the same from its outputs.
Quantifiers \code!s!, \code!t!; generated converse \cmplr{binder_definition_converse}.

\cmplhead{qf_dd}{[P2] DD: the splitter is a pink dot with plain mixed green dots.}
\cmpleq{\begin{tikzpicture}[zxd]
\node[snode] (s) at (0,0.3) {$S$};
\draw[wire] (0,-0.5) node[below,dimlabel] {$mn$} -- (s);
\draw[wire] (s.north west) -- ++(0,0.7) node[above,dimlabel] {$m$}; \draw[wire] (s.north east) -- ++(0,0.7) node[above,dimlabel] {$n$};
\end{tikzpicture}}{}{\begin{tikzpicture}[zxd]
\node[cmplpd] (p) at (0,0) {}; \node[cmplmultt] (m) at (-0.6,0.7) {$n$};
\node[cmplbox] (z1) at (-0.6,1.5) {}; \node[cmplbox] (z2) at (0.6,1.5) {};
\draw[wire] (0,-0.6) node[below,dimlabel] {$mn$} -- (p); \draw[wire] (p) -- (m) -- (z1); \draw[wire] (p) -- (z2);
\draw[wire] (z1) -- (-0.6,2.2) node[above,dimlabel] {$m$}; \draw[wire] (z2) -- (0.6,2.2) node[above,dimlabel] {$n$};
\end{tikzpicture}}
Left: $S$ from $mn$ to $(m,n)$. Right: a pink dot of dimension $mn$ with one input and two outputs;
output $0$ through the transposed multiplier labeled $n$ (at $mn$) into a plain mixed green dot with
an input of dimension $mn$ and an output of dimension $m$; output $1$ into a plain mixed green dot
with ports $mn$ and $n$. $\sigma=(mn)^{(n-3)/2}$. Left arities: for an $S$ anchor whose output
dimensions resolve to concrete $m$, $n$, yield them. Right arities: for an X anchor, every
$(m,n)$ of \py!_divisor_pairs(_host_values)!. Quantifiers \code!m!, \code!n!; generated converse
\cmplr{qf_dd_converse}.

\cmplhead{qf_e}{[P2] E: a pink dot state through a multiplier into dimension 2.}
\cmpleq{\begin{tikzpicture}[zxd]
\node[cmplpd] (p) at (0,0) {}; \node[cmplmult] (m) at (0,0.7) {$e$}; \node[cmplbox] (z) at (0,1.45) {};
\draw[wire] (p) -- node[right,dimlabel] {$n$} (m) -- (z) -- (0,2.1) node[above,dimlabel] {$2$};
\end{tikzpicture}}{}{\begin{tikzpicture}[zxd]
\node[cmplpd] (p) at (0,0.6) {}; \draw[wire] (p) -- (0,1.6) node[above,dimlabel] {$2$};
\end{tikzpicture}}
Left: a pink dot state of dimension $n$ through the multiplier labeled $e$ (at $n$) into a plain
mixed green dot with ports $n$ (input) and $2$ (output). Right: a pink dot state of dimension $2$.
$\sigma=n^{(3-e)/2}$. Left arities (\py!_e_arities!): for an X anchor with no inputs at a concrete
dimension $n$ whose output reaches a Z spider with $L$ ports, yield $(n,e)$ for
$e=1,\dots,\min(n,L+1)-1$ (the multiplier's top dot has $e+1$ ports, so $e\le L$). Right arities:
for an X anchor of dimension $2$, yield $(n,e)$ for every $n$ in \py!_host_values! and
$e=1,\dots,n-1$. Quantifiers \code!n!, \code!e!; generated converse \cmplr{qf_e_converse}.

\cmplhead{qf_dzx}{[P2] DZX: a pink dot of dimension $mn$ seen through dimension $m$.}
\cmpleq{\begin{tikzpicture}[zxd]
\node[cmplpd] (c) at (0,0.6) {};
\node[cmplmult] (mi) at (0,-0.2) {$n$}; \node[cmplbox] (zi) at (0,-0.95) {};
\node[cmplmultt] (mo) at (0,1.4) {$n$}; \node[cmplbox] (zo) at (0,2.15) {};
\draw[wire] (0,-1.6) node[below,dimlabel] {$m$} -- (zi) -- node[right,dimlabel] {$mn$} (mi) -- (c) -- (mo) -- (zo) -- (0,2.8) node[above,dimlabel] {$m$};
\end{tikzpicture}}{}{\begin{tikzpicture}[zxd]
\node[cmplpd] (c) at (0,0.6) {};
\draw[wire] (0,-0.4) node[below,dimlabel] {$m$} -- (c) -- (0,1.6) node[above,dimlabel] {$m$};
\end{tikzpicture}}
Shown for $p=q=1$. Left: a pink dot of dimension $mn$ with $p$ inputs and $q$ outputs; each input
$k$ is fed through a plain mixed green dot with ports $m$ (input, the terminal $i_k$) and $mn$
(output) and the multiplier labeled $n$ (at $mn$); each output passes through the transposed
multiplier labeled $n$ (at $mn$) into a plain mixed green dot with ports $mn$ and $m$ (output,
the terminal $o_k$). Right: a pink dot of dimension $m$ with $p$ inputs and $q$ outputs. With
$k=p+q$, $\sigma=m^{k-1}(mn)^{1-k(n-1)/2}$. Left arities: for an X anchor of concrete dimension
$D$, yield $(m,n,p,q)$ for every divisor pair with $mn=D$, $p$ and $q$ the anchor's input and output
counts. Right arities: for an X anchor of concrete dimension $m$, yield $(m,n,p,q)$ for every
$n\ge2$ in \py!_host_values!. Quantifiers \code!m!, \code!n!, \code!p!, \code!q!; promoted;
generated converse \cmplr{qf_dzx_converse}.

\cmplhead{qf_k0}{[P2] K0: a pink $K_{-j}$ state copied into two dimensions.}
\cmpleq{\begin{tikzpicture}[zxd]
\node[cmplpk] (p) at (0,0) {$K_{-j}$}; \node[cmplbox] (z) at (0,0.85) {};
\draw[wire] (p) -- node[right,dimlabel] {$m$} (z);
\draw[wire] (z) -- (-0.4,1.6) node[above,dimlabel] {$n_1$}; \draw[wire] (z) -- (0.4,1.6) node[above,dimlabel] {$n_2$};
\end{tikzpicture}}{[\,j\bmod m<\min(m,n_1,n_2)\,]\cdot}{\begin{tikzpicture}[zxd]
\node[cmplpk] (p1) at (-0.5,0.4) {$K_{-j}$}; \node[cmplpk] (p2) at (0.5,0.4) {$K_{-j}$};
\draw[wire] (p1) -- (-0.5,1.6) node[above,dimlabel] {$n_1$}; \draw[wire] (p2) -- (0.5,1.6) node[above,dimlabel] {$n_2$};
\end{tikzpicture}}
Left: a pink state of dimension $m$ with \py!Kj("j", sign=-1)! into a plain mixed green dot with
ports $m$ (input), $n_1$, $n_2$ (outputs). Right: pink $K_{-j}$ states of dimensions $n_1$ and
$n_2$; the right template carries the scalar $1$ when $(j\bmod m)<\min(m,n_1,n_2)$ and $0$
otherwise. $\sigma=m^{3/2}(n_1n_2)^{1/2}\cdot[\,j\bmod m<\min(m,n_1,n_2)\,]$. A forward
application may introduce the scalar $0$; the converse rejects a match whose matched template
scalar is $0$. Left arities: for an X anchor with no inputs at concrete dimension $m$, yield
$(m,n_1,n_2)$ for every $n_1\le n_2$ in \py!_host_values!. Right arities: for an X anchor with no
inputs at concrete dimension $n_1$, yield $(m,n_1,n_2)$ for every $m$ and every $n_2\ge n_1$ in
\py!_host_values!. Quantifiers \code!m!, \code!n1!, \code!n2!, \code!j!; promoted; generated
converse \cmplr{qf_k0_converse}.

\cmplhead{qf_bzw}{[P2] BZW: W nodes of dimensions $m$ and $n$ crossed into mixed Z boxes.}
\cmpleq{\begin{tikzpicture}[zxd]
\node[wnode] (wm) at (-0.5,0) {}; \node[wnode] (wn) at (0.5,0) {};
\node[cmplbox] (z1) at (-0.5,1.1) {$a$}; \node[cmplbox] (z2) at (0.5,1.1) {$a$};
\draw[wire] (wm) -- (z1); \draw[wire] (wm) -- (z2); \draw[wire] (wn) -- (z1); \draw[wire] (wn) -- (z2);
\draw[wire] (wm) -- (-0.5,-0.6) node[below,dimlabel] {$m$}; \draw[wire] (wn) -- (0.5,-0.6) node[below,dimlabel] {$n$};
\draw[wire] (z1) -- (-0.5,1.8) node[above,dimlabel] {$k$}; \draw[wire] (z2) -- (0.5,1.8) node[above,dimlabel] {$k$};
\end{tikzpicture}}{}{\begin{tikzpicture}[zxd]
\node[cmplbox] (z) at (0,0.2) {$a$}; \node[wnode] (w) at (0,1.0) {}; \draw[wire] (z) -- node[right,dimlabel] {$k$} (w);
\draw[wire] (z) -- (-0.4,-0.6) node[below,dimlabel] {$m$}; \draw[wire] (z) -- (0.4,-0.6) node[below,dimlabel] {$n$};
\draw[wire] (w) -- (-0.4,1.8) node[above,dimlabel] {$k$}; \draw[wire] (w) -- (0.4,1.8) node[above,dimlabel] {$k$};
\end{tikzpicture}}
Left: W nodes of dimensions $m$ and $n$ with two outputs each; output $c$ of each feeds the
mixed-dimensional Z box number $c$ with ports $m$, $n$ (inputs) and $k$ (output) and parameter
\py!MixedVar("a", min(m, n, k))! (both boxes bind the same $a$). Right: one such Z box into a W
node of dimension $k$ with two outputs. $\sigma=mnk$. Left arities: for a W anchor of concrete
dimension $m$, yield $(m,n,k)$ for every $n$ and $k$ in \py!_host_values!. Right arities: for a Z
anchor of concrete dimension $D$, every $(m,n,k)$ from \py!_host_values!$^3$ with $mnk=D$.
Quantifiers \code!m!, \code!n!, \code!k!; promoted; generated converse \cmplr{qf_bzw_converse}.

\cmplhead{qf_ta}{[P2] TA: the mixed trialgebra rule.}
\cmpleq{\begin{tikzpicture}[zxd]
\node[wnode] (w) at (-0.6,0) {}; \node[cmplpd] (p) at (0.6,-0.2) {};
\node[cmplmultt] (e1) at (0.2,0.55) {$e_1$}; \node[cmplmultt] (e2) at (1.0,0.55) {$e_2$};
\node[cmplbox] (z1) at (-0.4,1.35) {}; \node[cmplbox] (z2) at (0.6,1.35) {};
\draw[wire] (w) -- (z1); \draw[wire] (w) -- (z2); \draw[wire] (p) -- (e1) -- (z1); \draw[wire] (p) -- (e2) -- (z2);
\draw[wire] (w) -- (-0.6,-0.7) node[below,dimlabel] {$m$}; \draw[wire] (p) -- (0.6,-0.7) node[below,dimlabel] {$n$};
\draw[wire] (z1) -- (-0.4,2.0) node[above,dimlabel] {$k$}; \draw[wire] (z2) -- (0.6,2.0) node[above,dimlabel] {$k$};
\end{tikzpicture}}{}{\begin{tikzpicture}[zxd]
\node[wnode] (w) at (-0.6,0) {}; \node[wnode] (p) at (0.6,-0.2) {};
\node[cmplmultt] (e1) at (0.2,0.55) {$e_1$}; \node[cmplmultt] (e2) at (1.0,0.55) {$e_2$};
\node[cmplbox] (z1) at (-0.4,1.35) {}; \node[cmplbox] (z2) at (0.6,1.35) {};
\draw[wire] (w) -- (z1); \draw[wire] (w) -- (z2); \draw[wire] (p) -- (e1) -- (z1); \draw[wire] (p) -- (e2) -- (z2);
\draw[wire] (w) -- (-0.6,-0.7) node[below,dimlabel] {$m$}; \draw[wire] (p) -- (0.6,-0.7) node[below,dimlabel] {$n$};
\draw[wire] (z1) -- (-0.4,2.0) node[above,dimlabel] {$k$}; \draw[wire] (z2) -- (0.6,2.0) node[above,dimlabel] {$k$};
\end{tikzpicture}}
Shown for $i=2$. Left: a W node of dimension $m$ and a pink dot of dimension $n$, each with one
input and $i$ outputs; output $c$ of the pink dot passes through the transposed multiplier labeled
$e_c$ (at $n$); output $c$ of the W node and that multiplier feed a plain mixed green dot with ports
$m$, $n$ (inputs) and $k$ (output). Right: the same with the pink dot replaced by a W node of
dimension $n$. $\sigma=n^{-(i-1)/2}$. Both sides: for a W anchor of concrete dimension $m$ with $i$
outputs, yield $(m,n,k,i,e_0,\dots,e_{i-1})$ for every $n$, $k$ in \py!_host_values! and every
multiset $e_0\le\dots\le e_{i-1}$ from $1,\dots,n-1$. Quantifiers \code!m!, \code!n!, \code!k!,
\code!i!, \code!e_1..e_i!; generated converse \cmplr{qf_ta_converse}.

\cmplhead{qf_s1}{[P2] S1: mixed-dimensional Z boxes joined by wires fuse.}
\cmpleq{\begin{tikzpicture}[zxd]
\node[cmplbox] (a) at (0,0) {$a$}; \node[cmplbox] (b) at (0,1.4) {$b$};
\draw[wire] (a) to[bend left=40] node[left,dimlabel] {$d_1$} (b); \draw[wire] (a) to[bend right=40] node[right,dimlabel] {$d_2$} (b);
\draw[wire] (a) -- (0,-0.7); \draw[wire] (b) -- (0,2.1);
\end{tikzpicture}}{}{\begin{tikzpicture}[zxd]
\node[cmplbox] (c) at (0,0.7) {$ab'$}; \draw[wire] (c) -- (0,-0.7); \draw[wire] (c) -- (0,2.1);
\end{tikzpicture}}
The anchor and its partner are host mixed-dimensional Z boxes, recognized by
\py!mixed_legs(diagram, z)!. That function returns, for a Z spider \code!z!, the list of
(direction, dimension, external port) of its ports, or \py!None! unless every port of \code!z! is
wired as in \cref{def:completeness-mixedz}: an input port to the output of a $B$ whose input $0$
carries a phaseless X spider with exactly one port and whose input $1$ is the external port of
concrete dimension $d_i$; an output port to the input of an $S$ whose output $0$ carries such an X
spider and whose output $1$ is the external port; and the product of the $d_i$ equals the
dimension of \code!z! (read from its first port).

Left arities (\py!_s1m_arities!): for an anchor with legs $A$, every other node with legs $B$
registers its external ports; for each port of $A$ whose external port is wired to a registered
external port of some $B$, record the pair of indices; for each partner $B$ in increasing node id,
yield the arity map with \code!na!, \code!nb! (leg counts), \code!nc! (connections),
\code!a{k}d!, \code!a{k}o! (dimension, and $1$ for an output) for each leg of $A$, the same for
$B$ with prefix \code!b!, and \code!c{k}a!, \code!c{k}b! for each connection. The left template is
the mixed Z box $a$ over the legs of $A$ with \py!MixedVar("a", low_a)! and the mixed Z box $b$
over the legs of $B$ with \py!MixedVar("b", low_b)!, the connected external ports wired, and the
other ports terminals \code!a{k}! and \code!b{k}!. With $\ell=\min(\ell_a,\ell_b)$ and $\ell'$ the
least dimension among the free ports, the right template is the mixed Z box over the free ports
(in the order: those of $A$, then those of $B$) with the derived parameter
$ab'=(a_1b_1,\dots,a_{\ell-1}b_{\ell-1},0,\dots,0)$ over dimension $\ell'$, spread to the box's
dimension. When no port is free, the right template is empty and carries the scalar
$\sum_{k=0}^{\ell-1}a_kb_k$ (with $a_0=b_0=1$). The engine scalar is
\[
\sigma=\frac{\prod_{i\in\text{free}}(D_f/d_i)^{-1/2}}{\prod_{i\in A}(D_A/d_i)^{-1/2}\prod_{i\in B}(D_B/d_i)^{-1/2}}\cdot s_R,
\]
with $D_X$ the product of the port dimensions of $X$ and the free product read as $1$ when nothing
is free. Promoted; quantifiers \code!legs!, \code!dimensions!. The rule has no converse: the
right side binds neither $a$ nor $b$, so \py!_determined! rejects it, and no explicit converse is
written.


\section{Definitional and Structural Rules}\label{sec:completeness-definitional}

Five rules of \py!AXIOMS! translate engine generators and parameters into the papers' notation;
\cmplr{weight_to_phase}, an explicit converse, reverses one of them. With \cmplr{binder_definition}
(\cref{sec:completeness-mixed}) they are the rules of the reachability step of
\cref{sec:completeness-transfer}.

\cmplhead{x_spider_definition}{The X spider is the Z spider with $F$ on each output and $F^3$ on each input.}
\cmpleq{\begin{tikzpicture}[zxd]
\node[xs] (x) at (0,0.7) {$a$};
\draw[wire] (-0.4,-0.3) -- (x); \draw[wire] (0.4,-0.3) -- (x); \draw[wire] (x) -- (-0.4,1.7); \draw[wire] (x) -- (0.4,1.7);
\node[dimlabel] at (0,-0.25) {$n$}; \node[dimlabel] at (0,1.65) {$m$};
\end{tikzpicture}}{}{\begin{tikzpicture}[zxd]
\node[zs] (p) at (0,0.7) {$a$};
\node[fbox] (a1) at (-0.5,-0.2) {$F^3$}; \node[fbox] (a2) at (0.5,-0.2) {$F^3$};
\node[fbox] (b1) at (-0.5,1.6) {$F$}; \node[fbox] (b2) at (0.5,1.6) {$F$};
\draw[wire] (-0.5,-0.8) -- (a1) -- (p); \draw[wire] (0.5,-0.8) -- (a2) -- (p);
\draw[wire] (p) -- (b1) -- (-0.5,2.2); \draw[wire] (p) -- (b2) -- (0.5,2.2);
\node[dimlabel] at (0,-0.2) {$\cdots n$}; \node[dimlabel] at (0,1.6) {$\cdots m$};
\end{tikzpicture}}
Left: an engine X spider (not pink) with $n$ inputs, $m$ outputs, and \py!Var("a")!. Right: a Z
spider with \py!Var("a")!, three $F$ before each input and one after each output. $\sigma=1$.
Left arities: for an X anchor, yield its $(n,m)$. Right arities: those of \cmplr{zxw_h1}.
Quantifiers \code!n!, \code!m!, \code!d!; generated converse \cmplr{x_spider_definition_converse}.

\cmplhead{triangle_definition}{[P1] (YT): the triangle is a W node with one output on a plain green effect.}
\cmpleq{\begin{tikzpicture}[zxd]
\node[tri] (t) at (0,0.4) {}; \draw[wire] (0,-0.4) -- (t) -- (0,1.2);
\end{tikzpicture}}{}{\begin{tikzpicture}[zxd]
\node[wnode] (w) at (0,0.2) {}; \node[zdot] (e) at (-0.45,0.85) {};
\draw[wire] (0,-0.4) -- (w); \draw[wire] (w) -- (e); \draw[wire] (w) -- (0.45,1.2);
\end{tikzpicture}}
Left: $T$ (one input, one output). Right: a W node with two outputs, output $0$ into a plain green
effect, output $1$ the terminal \code!out!. $\sigma=1$; fixed sides; generated converse
\cmplr{triangle_definition_converse}.

\cmplhead{triangle_inverse_definition}{[P1] (YT): the inverse triangle is a W node with one output on the box $-1$.}
\cmpleq{\begin{tikzpicture}[zxd]
\node[tri] (t) at (0,0.4) {}; \node[dimlabel] at (0.4,0.4) {$-1$}; \draw[wire] (0,-0.4) -- (t) -- (0,1.2);
\end{tikzpicture}}{}{\begin{tikzpicture}[zxd]
\node[wnode] (w) at (0,0.2) {}; \node[cmplbox] (e) at (-0.55,0.9) {$-1$};
\draw[wire] (0,-0.4) -- (w); \draw[wire] (w) -- (e); \draw[wire] (w) -- (0.45,1.2);
\end{tikzpicture}}
Left: $T_i$. Right: as for \cmplr{triangle_definition} with the effect a Z spider carrying
\py!MINUS_ONES!, the Z box $(-1,\dots,-1)$. $\sigma=1$; fixed sides; generated converse
\cmplr{triangle_inverse_definition_converse}.

\cmplhead{phase_to_weight}{[P1] green circle spider: the phase $\alpha$ is the weight $e^{i\alpha}$.}
Left: an open Z spider with no template ports and \py!Var("a")!. Right: an open Z spider \code!z!
with the derived \py!as_weights(a)! that receives the left node's ports
(\py!rest = {"z": "z"}!). The left side's arity enumerator yields one empty assignment for a Z
anchor. Condition \cmplr{carries_a_phase}: the bound vector has an explicit phase entry or a phase
rule. $\sigma=1$. \py!as_weights(pv, dim)! returns the weight vector with the same coefficients:
at a concrete dimension \py!from_weights({k: pv.coefficient(k)})! for $k=1,\dots,d-1$; at a
symbolic dimension, for a vector with neither entries nor explicit weights, the weight rule
$r_{pv}$ of \cref{sec:completeness-vectors}. It raises \py!RewriteDomainError! with
\code!"a phaseless spider has no phase to move"! for \py!None!,
\code!"the vector carries no phase"! for a vector with no entry and no rule, and
\code!"moving explicit entries needs a concrete dimension"! at a symbolic dimension with explicit
entries or weights. The rule has no generated converse (it is excluded by name); its converse is
\cmplr{weight_to_phase}.

\cmplhead{weight_to_phase}{The converse of \cmplr{phase_to_weight}: weights that are roots of unity become phases.}
Left: an open Z spider with \py!Var("w")!; right: an open Z spider with the derived
\py!_as_phases(w)!. \py!concrete_only!; $\sigma=1$. \py!_as_phases! requires a vector of weights
only (no phase entries, no phase rule, \py!is_weighted!), else \py!RewriteDomainError!
\code!"weight_to_phase reads a vector of weights only"!; it returns the phase vector whose entry at
$k$ is \py!_phase_of_weight(w_k)!. \py!_phase_of_weight(value)! returns the phase $t$ (in turns)
with $\mathit{value}=e^{2\pi i t}$:
\begin{enumerate}
\item If the expression is \py!exp(x)! and $x/(2\pi i)$ simplifies to a real number $t$, return $t$.
\item If $|\mathit{value}|-1$ does not simplify to $0$, raise \code!"{value} is not a unit complex number"!.
\item Let $t$ be \py!nsimplify(arg(value) / (2*pi))!; if $t$ is not rational, raise
  \code!"{value} is not a rational root of unity"!.
\item If $e^{2\pi i t}-\mathit{value}$ does not simplify to $0$, raise
  \code!"{value} is not e^(2 pi i {t})"!; otherwise return $t$.
\end{enumerate}

\cmplhead{scalar_box}{A Z box with no ports is a scalar.}
\cmpleq{\begin{tikzpicture}[zxd]
\node[cmplbox] (z) at (0,0) {$a$};
\end{tikzpicture}}{\bigl(1+a_1+\dots+a_{d-1}\bigr)\cdot}{\cmplempty}
Left: a Z spider with no ports and \py!Var("a")! (not open: a host spider with ports does not
match). Right: the empty template. The axiom scalar is \py!_box_value(dim, a)!:
\begin{itemize}
\item at a concrete dimension, $\sum_{k=0}^{d-1}c_k$ over the coefficients ($c_0=1$);
\item at a symbolic dimension: for \py!None!, $d$; for a vector with explicit entries or weights,
  \py!RewriteDomainError! \code!"a symbolic sum over explicit entries"!; otherwise, with $r$ the
  rule coefficient $r_{a}$, the value $1+(d-1)r$ when $r$ does not contain \py!INDEX!, else
  $1+\sum_{k=0}^{d-1}r(k)-r(0)$ (\py!Scalar.index_sum!).
\end{itemize}
$\sigma$ equals that value. \py!converse_free = {"a": plain vector}!; promoted. The generated
converse is replaced by the explicit \cmplr{scalar_box_converse}
(\cref{sec:completeness-converses}).

\section{Templates of the Older Rules}\label{sec:completeness-legacy}

The hand-written rules of \cref{ch:rules} predate the template machinery and have no converses.
\py!LEGACY! restates 27 of their equations as template axioms, read left to right; only their
converses are registered (\py!LEGACY_CONVERSES!, named \cmplr{<name>_converse}), since the forward
direction is the hand-written rule. Every node is an engine node (no pink normalization), so
$\sigma$ is the axiom scalar. A side whose arity enumerator yields nothing (\py!lambda d, n, dim: ()!)
never matches; the converse matches the other side. The descriptions of these templates cite
rule numbers 2.1 to 2.28 of a retired document; the numbers carry no behavior.

\begin{table}[H]
\centering\small
\footnotesize
\begin{tabularx}{\linewidth}{@{}>{\raggedright\arraybackslash}p{0.40\linewidth}>{\raggedright\arraybackslash}X>{\raggedright\arraybackslash}p{0.10\linewidth}@{}}
\toprule
Template name & Left $=$ $\sigma\cdot$ Right & $\sigma$ \\
\midrule
\cmplr{fourier_cancellation} & four $F$ in series $=$ bare wire & $1$ \\
\cmplr{triangle_inverse_cancellation} & $T$ then $T_i$ $=$ bare wire & $1$ \\
\cmplr{triangle_inverse_cancellation_ti_t} & $T_i$ then $T$ $=$ bare wire & $1$ \\
\cmplr{w_identity} & W node with one output $=$ bare wire & $1$ \\
\cmplr{zx_cap} & plain Z state into plain X effect $=$ empty & $\sqrt d$ \\
\cmplr{zx_cap_swapped} & plain X state into plain Z effect $=$ empty & $\sqrt d$ \\
\cmplr{fourier_state_color_change} & plain Z state then $F$ $=$ plain X state & $1$ \\
\cmplr{fourier_effect_color_change} & $F$ then plain X effect $=$ plain Z effect & $1$ \\
\cmplr{fourier_state_color_change_swapped} & plain X state then $F$ $=$ plain Z state & $1$ \\
\cmplr{fourier_effect_color_change_swapped} & $F$ then plain Z effect $=$ plain X effect & $1$ \\
\cmplr{triangle_zero_state_t} & plain X state then $T$ $=$ plain X state & $1$ \\
\cmplr{triangle_zero_state_ti} & plain X state then $T_i$ $=$ plain X state & $1$ \\
\cmplr{triangle_zero_effect} & $T$ then plain X effect $=$ plain Z effect & $\sqrt d$ \\
\cmplr{w_z_effect} & W node ($m$ outputs), output $0$ on a plain Z effect $=$ $T$ then a W node with $m-1$ outputs & $1$ \\
\cmplr{w_zero_effect} & W node ($m$ outputs), output $0$ on a plain X effect $=$ a W node with $m-1$ outputs & $\sqrt d$ \\
\cmplr{hopf} & open Z spider $a$ (two outputs) and open X spider $b$ (two inputs) joined by a wire and by a wire through $F^2$ $=$ the two spiders disconnected & $1/d$ \\
\cmplr{hopf_swapped} & the same with the X spider below & $1/d$ \\
\cmplr{bialgebra_swapped} & plain Z (2 in, 1 out) into plain X (1 in, 2 out) $=$ two plain X (1 in, 2 out) crossed into two plain Z (2 in, 1 out) & $\sqrt d$ \\
\cmplr{wz_bialgebra} & plain Z ($n$ in, 1 out) into W ($m$ outputs) $=$ $n$ W nodes ($m$ outputs) crossed into $m$ plain Z ($n$ in, 1 out) & $1$ \\
\cmplr{state_copy} & plain X state into plain Z (1 in, $n$ out) $=$ $n$ plain X states & $d^{(1-n)/2}$ \\
\cmplr{state_copy_swapped} & plain Z state into plain X (1 in, $n$ out) $=$ $n$ plain Z states & $d^{(1-n)/2}$ \\
\cmplr{w_zero_copy} & plain X state into W ($n$ outputs) $=$ $n$ plain X states & $d^{(1-n)/2}$ \\
\cmplr{connective_inverse} & $S$ ($st\to s,t$) into $B$ $=$ bare wire & $1$ \\
\cmplr{connective_states_states} & plain Z states of dimensions $s$, $t$ into $B$ $=$ plain Z state of dimension $st$ & $1$ \\
\cmplr{connective_states_effects} & $S$ into plain Z effects of dimensions $s$, $t$ $=$ plain Z effect of dimension $st$ & $1$ \\
\cmplr{connective_states_swapped_states} & the same with plain X states & $1$ \\
\cmplr{connective_states_swapped_effects} & the same with plain X effects & $1$ \\
\bottomrule
\end{tabularx}
\caption{The 27 templates of \py!LEGACY!, in order.}
\label{tab:completeness-legacy}
\end{table}

\begin{center}
\begin{tikzpicture}[zxd]
\node[zdot] (z) at (0,0) {}; \node[xdot] (x) at (0,0.7) {}; \draw[wire] (z) -- (x);
\node at (0.6,0.35) {$=$}; \node at (1.3,0.35) {$\sqrt d$};
\begin{scope}[xshift=3cm]
\node[xs] (x) at (0,0) {$b$}; \node[zs] (z) at (0,1.4) {$a$}; \node[fbox] (f1) at (0.6,0.45) {$F$}; \node[fbox] (f2) at (0.6,0.95) {$F$};
\draw[wire] (z) to[bend right=30] (x); \draw[wire] (x) to[out=20,in=-90] (f1) -- (f2) to[out=90,in=-20] (z);
\draw[wire] (x) -- (-0.3,-0.6); \draw[wire] (z) -- (-0.3,2.0);
\node at (1.4,0.7) {$=$}; \node at (2.0,0.7) {$\frac1d$};
\node[xs] (x2) at (2.7,0.2) {$b$}; \node[zs] (z2) at (2.7,1.2) {$a$};
\draw[wire] (x2) -- (2.7,-0.6); \draw[wire] (z2) -- (2.7,2.0);
\end{scope}
\begin{scope}[xshift=8.2cm]
\node[xdot] (s) at (0,0) {}; \node[zdot] (g) at (0,0.7) {}; \draw[wire] (s) -- (g);
\draw[wire] (g) -- (-0.4,1.4); \draw[wire] (g) -- (0.4,1.4); \node[dimlabel] at (0,1.35) {$n$};
\node at (0.8,0.6) {$=$}; \node at (1.6,0.6) {$d^{\frac{1-n}{2}}$};
\node[xdot] (s1) at (2.3,0.5) {}; \node[xdot] (s2) at (2.9,0.5) {};
\draw[wire] (s1) -- (2.3,1.4); \draw[wire] (s2) -- (2.9,1.4); \node[dimlabel] at (2.6,1.0) {$\cdots$};
\end{scope}
\end{tikzpicture}
\end{center}
\noindent The pictures show \cmplr{zx_cap}, \cmplr{hopf} (the open spiders keep their other ports,
drawn as one wire each), and \cmplr{state_copy}; the full-red nodes are engine X spiders.

The arity enumerators of the converses are as follows. \cmplr{w_z_effect} and
\cmplr{w_zero_effect}: for a W anchor, $m=\text{outputs}+1$; for a $T$ anchor, $m=$ the output count
of an adjacent W node plus $1$, for each adjacent W node. \cmplr{hopf} and \cmplr{hopf_swapped}:
one empty assignment at a Z anchor (the right side is two open spiders with no ports, receiving
the rest ports: \py!rest = {"z": "z2", "x": "x2"}! and back). \cmplr{wz_bialgebra}: for a W anchor,
$n$ ranges over the distinct input counts of adjacent Z spiders and $m$ is the anchor's output
count. \cmplr{state_copy}, \cmplr{state_copy_swapped}, \cmplr{w_zero_copy}: for a phaseless state
anchor of the copied generator (no inputs, one output), $n=2,\dots,s$ with $s$ the number of such
phaseless states in the host. \cmplr{connective_inverse} (right side a bare wire): at a concrete
dimension $D$, every factor pair $(s,t)$ with $2\le s<D$, $s\mid D$, and $t=D/s\ge2$, enumerated
per wire of that dimension. \cmplr{connective_states_*}: for an anchor of the right generator
with exactly one port, the factor pairs of its dimension. The connective templates are
\py!concrete_only!.

\cmplhead{x_fusion_converse}{Spider fusion of X spiders read right to left.}
Left: an X spider \code!x! with $n_i$ inputs, $n_o$ outputs, and \py!Var("c")!. Right: X spiders
\code!xa! (parameter \py!Var("a")!) and \code!xb! (\py!Var("b")!); input $k$ goes to \code!xa! when
bit $k$ of the mask $m_i$ is set and to \code!xb! otherwise, output $k$ likewise by the mask
$m_o$; an extra output of \code!xa! feeds an extra input of \code!xb!. Arities: for an X anchor,
every $(n_i,n_o,m_i,m_o)$ with $0\le m_i<2^{n_i}$ and $0\le m_o<2^{n_o}$. Condition
\cmplr{fuses_back}: $ab=c$ (\py!_fuses_to!). Free variables: $a$ defaults to $c$, $b$ to the plain
vector. \py!solve = _solve_product!: given $a$ and $c$, $b_k=c_k/a_k$ for $k=1,\dots,d-1$, or
\py!None! when some $a_k$ is $0$ (by \py!_zero!). \py!candidates = _swapped_default!: the single
assignment $a=$ plain, $b=c$. $\sigma=1$. It belongs to \py!LEGACY_CONVERSES!.

\section{Converses}\label{sec:completeness-converses}

Every equation of \py!AXIOMS! except \cmplr{qf_s1} is available in both directions.

\begin{definition}[Determined axiom]\label{def:completeness-determined}
\py!_determined(axiom)! reads both templates at $d=3$ and at the first arity sample of
\py!_SAMPLES! for which both build without \py!KeyError!, \py!IndexError!, \py!ValueError!, or
\py!RewriteDomainError!. The samples are, in order,
$\{k{=}1,n{=}2,m{=}2,p{=}0,q{=}1,r{=}1,e{=}1,e0{=}1\}$, $\{k{=}0,n{=}2,m{=}2\}$, $\{n{=}1,e0{=}1\}$,
$\{s{=}2,t{=}3\}$, $\{m{=}2,n{=}3,k{=}2,i{=}1,e0{=}1\}$, $\{m{=}2,n1{=}2,n2{=}3\}$. The axiom is
determined when both templates exist, the right template has a node, a bare wire, or
\py!converse_free! is non-empty, and the variables read by the left template's parameters are
among those bound by the right template's parameters together with the names of
\py!converse_free!. No sample building both sides gives false.
\end{definition}

\py!CONVERSES! is the concatenation of three groups:
\begin{enumerate}
\item the generated converses \py!axiom.converse(name + "_converse")! of every axiom of \py!AXIOMS!
  that is determined, whose converse name is not among the explicit converses, and that is not
  \cmplr{phase_to_weight}: the 36 rules \cmplr{zxw_s2}, \cmplr{zxw_s2_pink}, \cmplr{zxw_ept},
  \cmplr{zxw_b2}, \cmplr{zxw_k0}, \cmplr{zxw_k1}, \cmplr{zxw_zer}, \cmplr{zxw_p1},
  \cmplr{zxw_p1_multiplier}, \cmplr{zxw_h1}, \cmplr{zxw_pcy}, \cmplr{zxw_sym}, \cmplr{zxw_aso},
  \cmplr{zxw_bzw}, \cmplr{zxw_ww}, \cmplr{zxw_bs0}, \cmplr{zxw_ta}, \cmplr{zxw_hd},
  \cmplr{zxw_vw}, \cmplr{zxw_kz}, \cmplr{qf_wf}, \cmplr{qf_w_sym}, \cmplr{qf_z_flex},
  \cmplr{qf_mw}, \cmplr{qf_ww}, \cmplr{qf_hd}, \cmplr{binder_definition}, \cmplr{qf_dd},
  \cmplr{qf_e}, \cmplr{qf_dzx}, \cmplr{qf_k0}, \cmplr{qf_bzw}, \cmplr{qf_ta},
  \cmplr{x_spider_definition}, \cmplr{triangle_definition}, and
  \cmplr{triangle_inverse_definition}, each followed by \cmplr{_converse};
\item \py!EXPLICIT_CONVERSES! (11 rules), whose matched side carries a derived parameter;
\item \py!LEGACY_CONVERSES! (28 rules): the 27 converses of \py!LEGACY! and
  \cmplr{x_fusion_converse}.
\end{enumerate}
The explicit converses are listed below. Each is \py!concrete_only! except where stated; the
scalar of each follows \cref{eq:completeness-engine-scalar}.

\begin{description}[style=nextline,leftmargin=1.2em]
\item[\cmplr{scalar_box_converse}] The converse of \cmplr{scalar_box} with the condition
  \cmplr{box_value_nonzero} added (\py!dataclasses.replace! then \py!converse!):
  inserts a Z box with no ports and the default plain vector (value $d$), times $1/v$ with $v$ its
  value; condition \cmplr{box_value_nonzero}: $v\neq0$ by \py!_nonzero!. Not \py!concrete_only!.
\item[\cmplr{zxw_d1_converse}] Two $F$ then a Z spider \py!Var("b")! become a Z spider
  \py!reversed_vector(b)! then two $F$.
\item[\cmplr{zxw_k2_converse}] A Z spider \py!Var("c")! into a pink $K_j$ becomes a pink $K_j$ into
  a Z spider with \py!_unshift(c, j)!, $a_k=c_{k+j}/c_j$, and the right template carries the scalar
  $c_j$: $c\,;K_j=c_j\cdot K_j\,;a$. Condition \cmplr{c_j_nonzero}; \py!_unshift! raises
  \code!"k_j(a) = c needs c_j != 0"! otherwise.
\item[\cmplr{zxw_s1_converse}] Unfusion of a Z box: left, a Z spider with no inputs, $n$ outputs, and
  \py!Var("c")!; right, Z spiders $a$ (no inputs) and $b$ joined by $k$ wires, port $p$ of the
  original going to $a$ when bit $p$ of the mask is set and to $b$ otherwise. Arities: for a Z
  anchor with $n$ ports, every mask $0\le\mathit{mask}<2^n$ and $k\in\{1,2\}$. Condition
  \cmplr{fuses_back} ($ab=c$); free variables $a=c$ and $b=$ plain by default;
  \py!solve = _solve_product!; \py!candidates = _single_entry_splits!: the assignment
  $(a,b)=(\text{plain},c)$ and, at a concrete dimension when at least two indices $k$ of $c$ have
  a coefficient other than $1$, for each such $k$ the two assignments $(c^{(k)},c^{\neg k})$ and
  $(c^{\neg k},c^{(k)})$, where $c^{(k)}$ carries the weight $c_k$ at $k$ alone and $c^{\neg k}$
  the weights at the other such indices. Not \py!concrete_only!.
\item[\cmplr{zxw_ad_converse}] A Z effect \py!Var("c")! becomes a W node with effects $a$ and $b$.
  Condition \cmplr{adds_back} ($a+b=c$ by \py!_same_coefficients!); free variables $a=c$,
  $b=$ \py!zero_vector! by default; \py!solve = _solve_sum! ($b_k=c_k-a_k$).
\item[\cmplr{zxw_bsj_converse}] A Z state matching \py!UnitIndex("j", sign=-1)! becomes the left
  side of \cmplr{zxw_bsj} with the bound $j$.
\item[\cmplr{zxw_zv_converse}] A Z state \py!Var("c")! becomes a Z state $a$ into V. Condition
  \cmplr{constant_vector}: every coefficient of $c$ equals $c_1$ and $c$ equals the constant vector
  of $a_{d-1}$; free variable $a$, default \py!last_vector(c_1)!.
\item[\cmplr{zxw_va_converse}] V then a Z spider \py!Var("a")! becomes the left side of
  \cmplr{zxw_va} with the boxes labelled $a_k$ derived from $a$ (\py!_box_entry(k)!).
\item[\cmplr{qf_va_converse}] The right side of \cmplr{qf_va} with \py!Var("a")! becomes its left
  side with derived boxes.
\item[\cmplr{qf_d1_converse}] A Z spider \py!Var("b")! with two $F$ after each output becomes a Z
  spider \py!reversed_vector(b)! with two $F$ before each input; arities as \cmplr{qf_d1}.
\item[\cmplr{weight_to_phase}] See \cref{sec:completeness-definitional}.
\end{description}

A converse whose matched side is empty or a bare wire inserts its built side: once per wire for a
bare wire (plus once per suitable boundary port when the built side takes no arities), and once
per distinct node dimension of the host for an empty side (\cref{ch:rewriting}).

\section{The Registry}\label{sec:completeness-registry}

\begin{signature}
\py!AXIOMS: tuple[TemplateAxiom, ...]! --- the 48 axioms and definitional rules, in the order of
\cref{sec:completeness-p1,sec:completeness-p2,sec:completeness-definitional}: the 27 rules of [P1]
(\cmplr{zxw_s1} to \cmplr{zxw_kz}), \cmplr{qf_wf}, \cmplr{qf_w_sym}, \cmplr{qf_z_flex},
\cmplr{qf_mw}, \cmplr{qf_d1}, \cmplr{qf_ww}, \cmplr{qf_va}, \cmplr{qf_hd},
\cmplr{binder_definition}, \cmplr{qf_dd}, \cmplr{qf_e}, \cmplr{qf_dzx}, \cmplr{qf_k0},
\cmplr{qf_bzw}, \cmplr{qf_ta}, \cmplr{qf_s1}, \cmplr{x_spider_definition},
\cmplr{triangle_definition}, \cmplr{triangle_inverse_definition}, \cmplr{phase_to_weight},
\cmplr{scalar_box}.

\py!CONVERSES: tuple[TemplateAxiom, ...]! --- the 75 converses of \cref{sec:completeness-converses}.

\py!LEGACY!, \py!EXPLICIT_CONVERSES!, \py!LEGACY_CONVERSES! --- the groups above.

\py!by_name() -> Mapping[str, TemplateAxiom]! --- a new dict of every axiom and converse by name
(123 entries; names are unique).

\py!__all__ = ["AXIOMS", "CONVERSES", "by_name"]!.
\end{signature}

\py!TEMPLATE_RULES! of \module{archytaszx.rewrite.rules_library} applies \py!template_rule! to
\py!(*AXIOMS, *CONVERSES)!; the registry \py!RULES! holds these 123 rules after the 28
hand-written rules, 151 rules in all. The forward templates of \py!LEGACY! are not registered.
\module{archytaszx.rewrite.rules_discard} imports \py!TB!, \py!_legs!, \py!_partners!,
\py!fixed_side!, \py!i!, and \py!o! from this module.

\begin{table}[H]
\centering\footnotesize
\begin{tabularx}{\linewidth}{@{}>{\raggedright\arraybackslash}p{0.31\linewidth}>{\raggedright\arraybackslash}p{0.12\linewidth}>{\raggedright\arraybackslash}X>{\raggedright\arraybackslash}p{0.17\linewidth}@{}}
\toprule
Rule & Source & Engine scalar $\sigma$ & Converse \\
\midrule
\cmplr{zxw_s1} & [P1] S1 & $1$ & explicit \\
\cmplr{zxw_s2}, \cmplr{zxw_s2_pink} & [P1] S2 & $1$ & generated \\
\cmplr{zxw_d1} & [P1] D1 & $1$ & explicit \\
\cmplr{zxw_ept} & [P1] Ept & $\sqrt d$ & generated \\
\cmplr{zxw_b2} & [P1] B2 & $d^{-1/2}$ & generated \\
\cmplr{zxw_k0} & [P1] K0 & $d^{-1/2}$ & generated \\
\cmplr{zxw_k1} & [P1] K1 & $1$ & generated \\
\cmplr{zxw_k2} & [P1] K2 & $a_{-j}$ & explicit \\
\cmplr{zxw_zer} & [P1] Zer & $d^{-1/2}$ & generated \\
\cmplr{zxw_p1} & [P1] P1, (Du) & $1$ & generated \\
\cmplr{zxw_p1_multiplier} & [P1] P1 & $d^{(d-2)/2}$ & generated \\
\cmplr{zxw_h1} & [P1] H1 & $1$ & generated \\
\cmplr{zxw_pcy}, \cmplr{zxw_sym}, \cmplr{zxw_aso}, \cmplr{zxw_bzw}, \cmplr{zxw_ww} & [P1] Pcy, Sym, Aso, BZW, WW & $1$ & generated \\
\cmplr{zxw_ad} & [P1] AD & $1$ & explicit \\
\cmplr{zxw_bs0} & [P1] Bs0 & $d^{-1/2}$ & generated \\
\cmplr{zxw_bsj} & [P1] Bsj & $\sqrt d$ & explicit \\
\cmplr{zxw_ta} & [P1] TA & $d^{-1/2}$ & generated \\
\cmplr{zxw_hd} & [P1] HD & $d^{1/2-(d-2)(d^2+2d+3)/4}$ & generated \\
\cmplr{zxw_zv} & [P1] ZV & $d^{-(d/2-1)-(d-1)(d-2)/4}$ & explicit \\
\cmplr{zxw_va} & [P1] VA & $1$ & explicit \\
\cmplr{zxw_vw} & [P1] VW & $d^{d/2-1+(d-1)(d-2)/4}$ & generated \\
\cmplr{zxw_kz} & [P1] KZ & $d^{1-n-\sum(e_k-1)/2}$ & generated \\
\cmplr{qf_wf}, \cmplr{qf_w_sym}, \cmplr{qf_z_flex} & [P2] WF, (2.10), (2.9) & $1$ & generated \\
\cmplr{qf_mw} & [P2] MW & $d^{(e-1)/2}$ & generated \\
\cmplr{qf_d1} & [P2] D1 & $1$ & explicit \\
\cmplr{qf_ww} & [P2] WW & $1$ & generated \\
\cmplr{qf_va} & [P2] VA & $d^{-(d-2)(d+1)/4}$ & explicit \\
\cmplr{qf_hd} & [P2] HD & $(-1)^{d+1}d^{1/2-(d-2)(d+3)/4}$ & generated \\
\cmplr{binder_definition} & [P2] (2.8) & $1$ & generated \\
\cmplr{qf_dd} & [P2] DD & $(mn)^{(n-3)/2}$ & generated \\
\cmplr{qf_e} & [P2] E & $n^{(3-e)/2}$ & generated \\
\cmplr{qf_dzx} & [P2] DZX & $m^{p+q-1}(mn)^{1-(p+q)(n-1)/2}$ & generated \\
\cmplr{qf_k0} & [P2] K0 & $m^{3/2}(n_1n_2)^{1/2}[j\bmod m<\min(m,n_1,n_2)]$ & generated \\
\cmplr{qf_bzw} & [P2] BZW & $mnk$ & generated \\
\cmplr{qf_ta} & [P2] TA & $n^{-(i-1)/2}$ & generated \\
\cmplr{qf_s1} & [P2] S1 & see the rule & none \\
\cmplr{x_spider_definition} & definition & $1$ & generated \\
\cmplr{triangle_definition}, \cmplr{triangle_inverse_definition} & [P1] (YT) & $1$ & generated \\
\cmplr{phase_to_weight} & definition & $1$ & \cmplr{weight_to_phase} \\
\cmplr{scalar_box} & definition & $1+\sum_k a_k$ & explicit \\
\bottomrule
\end{tabularx}
\caption{The 48 rules of \py!AXIOMS! with their engine scalars.}
\label{tab:completeness-summary}
\end{table}

\section{The Transfer of Completeness}\label{sec:completeness-transfer}

The module does not state a theorem; its rules are chosen so that the following argument holds
for the registered rule set. A diagram is \term{closed} when it has no \bbox{} and its parameter
environment binds every free symbol; substituting the binding gives a diagram whose dimensions,
parameters, and scalar are values. A closed diagram is \term{single-dimensional} when every port
has one dimension and it has no $B$ or $S$ node. An \term{engine step} applies a registered rule at
a match its pattern verifies, with its free values set, or identifies two isomorphic diagrams;
it is forward when it replaces the rule's left side by its right side.

\begin{enumerate}
\item \emph{Translation.} The translation $\tau$ of \cref{tab:completeness-translation} acts on
  generators and wires separately and preserves the denotation, scalars included.
\item \emph{Reachability.} For a closed diagram $\mathcal D$, replace each X spider, $T$, $T_i$,
  phase vector, $B$, and scalar factor by its paper definition; call the result $\pi(\mathcal D)$,
  a diagram of $\mathrm{ZXW}_d$ when $\mathcal D$ is single-dimensional of dimension $d$ and of
  $\mathrm{ZXW}_{\mathrm{qf}}$ otherwise. Each replaced item is the left side of
  \cmplr{x_spider_definition}, \cmplr{triangle_definition}, \cmplr{triangle_inverse_definition},
  \cmplr{phase_to_weight}, \cmplr{binder_definition}, or \cmplr{scalar_box}, whose right side is
  $\tau$ of the definition; one forward step per occurrence rewrites $\mathcal D$ to
  $\tau(\pi(\mathcal D))$.
\item \emph{Lifting.} Every axiom $L=R$ of either calculus is a registered rule through $\tau$
  (\cref{tab:completeness-summary}), and every such rule except \cmplr{qf_s1} has a converse, so a
  paper step $C[L]\to C[R]$ or $C[R]\to C[L]$ becomes forward engine steps; the step from the right
  to the left side of the mixed S1 is one reverse step of \cmplr{qf_s1}. Matching is up to the
  symmetries the papers assume: all ports of a Z spider are interchangeable (flexsymmetry, with
  \cmplr{qf_z_flex}), the outputs of a W node (with \cmplr{zxw_sym} and \cmplr{qf_w_sym}), and an X
  spider's inputs and its outputs among themselves. Notations used in a form the pattern does not
  use ($H^\dagger$, $D$, pink spiders, scalar diagrams) are rewritten by the definitional rules and
  \cmplr{zxw_p1}; identities a paper step inserts on a wire are the converses of \cmplr{zxw_s2} and
  \cmplr{zxw_s2_pink}; converses with free variables are applied with the values of the paper's
  instance through \azx!with NAME=VALUE!, which the pattern verifies.
\item \emph{Completeness.} For closed $\mathcal D_1$, $\mathcal D_2$ with equal denotations,
  $\pi(\mathcal D_1)$ and $\pi(\mathcal D_2)$ are equal paper diagrams, so the completeness theorem
  of [P1] or [P2] gives a paper derivation between them; lifting each step and joining it with the
  two reachability derivations (the second read backward through the definitional converses) gives
  an engine derivation from $\mathcal D_1$ to $\mathcal D_2$. For single-dimensional diagrams no
  step involves \cmplr{qf_s1}, so every step is forward.
\end{enumerate}

\section{Exhaustive Search}\label{sec:completeness-exhaustive}

\module{archytaszx.rewrite.exhaustive} searches for an engine derivation. \term{Exhaustive search}
is a fair bidirectional breadth-first enumeration: from the start and from the goal, layer by layer,
every match of every rule of the rule set (automatic or not, converses included) is applied to
every diagram of the current layer.

\subsection{State space and deduplication}

The state of a side is a set of \term{reached} diagrams. A record
\py!_Reached(diagram, view, meet, results)! (frozen, slotted) holds a diagram, its comparison view
(\py!comparison_view!: spider ports in canonical order, empty parameter environment, simplified
scalar, the scalar $c\bar c$ for a mixed diagram; \cref{ch:rewriting}), the \term{meet view} on
which the two sides are compared (the comparison view, or \py!algebraic_view! in algebraic search),
and the tuple of \py!RewriteResult! records of the rule applications that produced it from the
side's root.

\begin{signature}
\py!class _Side(name: str, root: Diagram, values: Sequence[Scalar] | None = None)!\\
Fields: \py!name!, \py!values! (the algebraic values, or \py!None!); \py!seen: dict[str, list[_Reached]]!
keyed by \py!view_key(view)!; \py!met: dict[str, list[_Reached]]! keyed by \py!view_key(meet)!;
\py!frontier: list[_Reached]!; \py!depth: int!. Initially both maps and the frontier hold the root
record and \py!depth = 0!.
\begin{itemize}[leftmargin=1.2em]
\item \py!known(key, view)!: some record under \py!key! in \py!seen! has a view isomorphic to
  \py!view! (\py!views_isomorphic!, up to spider port order).
\item \py!lookup(key, meet)!: the first record under \py!key! in \py!met! whose meet view is
  isomorphic to \py!meet!, else \py!None!.
\item \py!expand(rules)!: one layer (\cref{alg:completeness-expand}); returns the list of
  (meet key, record) pairs added, in order.
\end{itemize}
\end{signature}

\py!view_key! is the canonical key of a comparison view up to spider port order; two diagrams with
equal keys are compared by isomorphism, so the key is a bucket and isomorphism decides.

\begin{algorithm}[Expanding one layer]\label{alg:completeness-expand}
For each record $r$ of the frontier, in order, and each rule of \py!rules!, in order:
\begin{enumerate}
\item Let $M$ be the matches that \py!find_matches! of the rule's pattern returns on
  \py!r.diagram!. If the pattern has
  \py!candidate_matches!, append, for each $m\in M$, its candidate matches. If the side has values and the
  pattern has \py!algebraic_matches!, append, for each $m$, those of \py!algebraic_matches(r.diagram, m, values)! not
  already in $M$. An exception of \py!_SKIPPED! during this step skips the rule for $r$.
\item For each match: apply the rule (\py!apply!); compute the comparison view and its key; skip the
  result when \py!known(key, view)!; compute the meet view and its key. An exception of
  \py!_SKIPPED! skips this match.
\item Store the new record (results extended by this application) in \py!seen! and \py!met!,
  append it to the next frontier and to the returned list.
\end{enumerate}
Then the frontier becomes the next frontier and \py!depth! increases by $1$ (also when the frontier
was empty).
\end{algorithm}

\py!_SKIPPED! is \py!APPLICATION_ERRORS! of \module{archytaszx.rewrite.egraph} (\py!RewriteError!,
\py!BangBoxError!, \py!DimensionError!, \py!GeneratorError!, \py!GraphError!, \py!PhaseError!,
\py!ScalarError!, \py!ValidateError!) plus \py!ValueError! and \py!ArithmeticError!; a diagram whose
construction or key raises one of them, for example a scalar beyond Python's integer conversion
limit, is skipped. Pruning is per side: a diagram reached on the start side is pruned only against
the start side's \py!seen!.

\subsection{The search}

\begin{signature}
\py!exhaustive_search(start: Diagram, goal: Diagram, *, rules: Sequence[Rule] | None = None, on_layer: Callable[[LayerReport], None] | None = None) -> ProofPath!\\
A derivation between \py!start! and \py!goal! over \py!rules! (default: every rule of \py!RULES!,
in registry order). Returns only on success; raises \py!ExhaustedError! when both sides are closed.
\end{signature}

\begin{algorithm}[Fair bidirectional search]\label{alg:completeness-search}
\begin{enumerate}
\item Build the sides \code!start! and \code!goal!. If a record of the start side's \py!met! is
  found by the goal side's \py!lookup!, return the path for that pair (no steps).
\item Repeat forever: for the start side and then the goal side, expand one layer; for each added
  (key, record), look it up in the other side; on a hit return the path with the start-side record
  first. After each side's expansion, call \py!on_layer! with a \py!LayerReport! and, if both
  frontiers are empty, raise \py!ExhaustedError!.
\end{enumerate}
\end{algorithm}

\py!_path(start, goal, left, right)! returns \py!ProofPath(start, goal, forward, backward)! with
\py!forward = (ProofStep("exhaustive", left.results),)! when \py!left.results! is non-empty (else
empty) and \py!backward! likewise from the goal-side record: the backward results rewrite the goal
toward the meeting diagram and are replayed in reverse by the proof certificate (\cref{ch:proofs}).

\begin{signature}
\py!@dataclass(frozen=True, slots=True) class LayerReport!: \py!side: str! (\code!"start"!,
\code!"goal"!, or \code!"start (bound K)"! / \code!"goal (bound K)"! in algebraic search),
\py!depth: int! (the side's depth after the layer), \py!reached: int! (records in \py!seen!),
\py!frontier: int! (records in the new frontier).

\py!class ExhaustedError(Exception)!, raised by \py!_exhausted()! with the message
\code!"both sides are closed under every rule and never meet"!.
\end{signature}

\paragraph{Fairness and termination.} Layer $k$ of the start side is expanded before layer $k$ of
the goal side, and both before layer $k+1$ of either. A derivation of length $n$ whose steps apply
rules at matches the patterns enumerate (a free variable at its default or at a value its rule's
candidate generator lists) is found once the start side reaches layer $n$, and once it reaches
layer $\lceil n/2\rceil$ and the goal side layer $\lfloor n/2\rfloor$ when the last
$\lfloor n/2\rfloor$ steps use rules with converses. There is no depth cap and no time limit: for
two diagrams that no enumerated derivation joins and whose reachable sets are infinite, the search
does not terminate. It ends with \py!ExhaustedError! only when both reachable sets are finite and
disjoint up to isomorphism.

\subsection{Algebraic search}

\begin{signature}
\py!algebraic_search(start, goal, *, rules=None, on_layer=None) -> ProofPath!\\
Iterative deepening over rounds $K=1,2,\dots$. Round $K$ takes the values \py!up_to_level(K)!
(every algebraic number of level at most $K$, each once, \module{archytaszx.algebra.algebraic}),
builds fresh sides with these values, checks the roots for a meeting, and then expands both sides
alternately, with lookups as in \cref{alg:completeness-search}, while the start side's depth is
below $\lceil K/2\rceil$ and some frontier is non-empty. Meet views are \py!algebraic_view!: the
comparison view with every spider parameter a weight vector and every value, scalar included, in
the canonical form of algebraic numbers (a rational, or \py!CRootOf(p, i)! with $p$ the primitive
minimal polynomial). It never raises \py!ExhaustedError!; without a derivation it runs forever.
\end{signature}

The level of an algebraic number is the height $\max(|p|,q)$ of a rational $p/q$; $n$ for a
primitive $n$-th root of unity with $n\ge3$; $1+\max$ of the heights of $a$ and $b$ for $a+bi$ with
$b\neq0$; and $D+H$ for a root of an irreducible integer polynomial of degree $D\ge2$ with largest
absolute coefficient $H$. Every algebraic number has a level and each level is finite, so every
assignment of algebraic values to the free weight-vector variables of a converse occurs in some
round. In a round, \py!algebraic_matches! sets the free variables that default to a phase vector to
every weight vector with entries among the values (all $|V|^{d-1}$ vectors per variable); with a
\py!solve! function and more than one free variable, the first variable ranges over the vectors and
the rest are solved, or enumerated when \py!solve! returns \py!None!. Building the start or goal side
raises when \py!algebraic_view! cannot canonicalize a value; \py!prove_exhaustive! checks this first
and raises \py!ProveGrammarError! \code!"... algebraic search needs concrete dimensions and algebraic parameters: ..."!.

\subsection{Use}

The command \azx!prove exhaustive [algebraic] [against] TARGET! calls \py!prove_exhaustive! of
\module{archytaszx.semantics.prove}, which runs the interface and oracle refutation rungs, calls
\py!exhaustive_search! (or \py!algebraic_search!), turns \py!ExhaustedError! into the status
\texttt{NOT\_FOUND} with reason \code!"exhaustive search: ..."!, and certifies and checks the path
(\cref{ch:proofs}). Diagrams with discards are first doubled.

\section{Coverage and Limits}\label{sec:completeness-limits}

\begin{itemize}
\item \emph{Covered.} Closed diagrams: concrete system size (no \bbox{}), concrete dimensions, and
  parameter values. By \cref{sec:completeness-transfer}, two such diagrams with equal denotations
  are joined by an engine derivation; \py!algebraic_search! enumerates every derivation whose free
  values are algebraic, so it finds one whenever such a derivation exists over algebraic
  parameters.
\item \emph{Not covered.} \bboxes{} and symbolic dimensions: neither paper proves completeness
  there, and the \py!concrete_only! rules do not match at a symbolic dimension.
\item \emph{Free values.} \py!exhaustive_search! applies converses with free variables only at
  their defaults and at the values of \py!candidates!; a derivation that needs another value, for
  example a split of a Z box not produced by \py!_single_entry_splits!, is not enumerated.
\item \emph{Arity enumerators.} A family instance is tried only when its rule's enumerator yields
  its arities; the mixed rules draw dimensions from the host's port dimensions, \cmplr{zxw_ww_converse}
  bounds $n$ by the host's node count, and the copy converses bound $n$ by the host's phaseless
  states. An instance outside these bounds is not reached in one step.
\item \begin{sloppypar}\emph{Converses.} \cmplr{z_fusion_prime_d}, \cmplr{identity_removal},
  \cmplr{bialgebra}, \cmplr{w_fusion}, \cmplr{connective_inverse_bind_first},
  \cmplr{port_box_unfusion}, and \cmplr{qf_s1}, and \cmplr{spider_fusion} for Z spiders, have no
  converse in the rule set; a derivation that needs them from right to left on the start side is
  found only from the goal side.\end{sloppypar}
\item \emph{Cost.} Every layer applies every rule at every match; the number of reached diagrams
  grows with the depth without bound. The search has no time or memory limit.
\end{itemize}

\section{Worked Examples}\label{sec:completeness-examples}

The directory \file{examples/completeness/} holds the lemmas of [P1] and [P2] at concrete
dimensions as scripts: \file{lemmas/p1_lemma07.azx} to \file{lemmas/p1_lemma48.azx} (with variants
such as \file{p1_lemma14b.azx}), \file{lemmas/p2_lemmaA03.azx} to \file{lemmas/p2_lemmaA80.azx} (with variants such as \file{p2_lemmaA40b.azx}), and
the explicit derivation \file{derivations/p1_lemma38_derivation.azx}. A paper pink spider is
written as an engine X spider with its normalization $d^{(m+n)/2-1}$ as scalar.

\begin{example}[Lemma 38 of {[P1]} by explicit steps]\label{ex:completeness-l38}
At $d=3$, a W node with one output on the pink dot effect is the identity. The script
\begin{lstlisting}[style=azx]
@rhs =
    z = Z(1 -> 1, literal 3)

w = W(1 -> 2, literal 3); e = X(1 -> 0, literal 3); wire w.out1 -- e.in0; scalar 3^(-1/2)
apply w_zero_effect
apply zxw_s2_converse at 0
apply w_identity
prove against @rhs
\end{lstlisting}
applies \cmplr{w_zero_effect} (scalar $\sqrt3$, cancelling $3^{-1/2}$, leaving a W node with one
output), inserts a plain green dot with \cmplr{zxw_s2_converse} (scalar $1$), and removes the W
node with \cmplr{w_identity}; the shell prints \code!step: w_zero_effect (consumed n0, n1, new n2, scalar sqrt(3))!
and the proof kernel re-checks the three steps.
\end{example}

\begin{example}[The same lemma by exhaustive search]\label{ex:completeness-l38x}
Replacing the three \azx!apply! lines and the last line by \azx!prove exhaustive against @rhs! makes
the search find a derivation itself; the shell prints
\code!prove: proved: proved exhaustively in 3 step(s), 2 forward and 1 backward; ...!: the start side
reached the meeting diagram at depth $2$ and the goal side at depth $1$.
\end{example}

\begin{example}[Lemma 7 of {[P1]}]\label{ex:completeness-l7}
\file{lemmas/p1_lemma07.azx} states, at $d=3$ and $j=1$, that the pink $K_1$ state into the Z box
effect $a=(a_1,a_2)$ is the scalar box $(0,a_2-1)$:
\begin{lstlisting}[style=azx]
@lhs =
    p = X(0 -> 1, literal 3) phase{k: k/3}; z = Z(1 -> 0, literal 3) weight{1: a1, 2: a2}; p -> z
    scalar 3^(-1/2)
@rhs =
    b = Z(0 -> 0, literal 3) weight[literal 3]{1: 0, 2: a2 - 1}
load @lhs
prove against @rhs
certificate kernel
\end{lstlisting}
\end{example}
